\documentclass[../main.tex]{subfiles}
\begin{document}
%==========================================TIÊU ĐỀ TẬP========================================
\begin{table}[h]
    \begin{adjustwidth}{-1cm}{}
        \begin{tabular}{>{\centering\arraybackslash}p{6cm}|>{\raggedright\arraybackslash}p{12.5cm}}
            \multirow{5}{6cm}
            {\\[-40pt]
                \flaregothic\fontsize{20pt}{20pt}\selectfont \centering
                \color{eptype}HISTOIRE\color{black}\\
                \freshbot\fontsize{72pt}{72pt}\selectfont
                \color{epnum}21\\[6pt]
                \cafeteria\fontsize{20pt}{20pt}\selectfont\color{epnum}(D-21)\color{black}
                \\[18pt]
                \color{black}
            }
             &
            {\fotskip\fontsize{18pt}{18pt}\selectfont \ruby{風}{かぜ}を\ruby{引}{ひ}いた\ruby{管}{かん}\ruby{理}{り}\ruby{者}{しゃ}}
            \\[6pt]
             & {\cafeteria\fontsize{18pt}{18pt}\selectfont When the Manager Falls Ill} \\[4pt]
             & {\vnmsans\fontsize{18pt}{18pt}\selectfont Cơn sốt bất ngờ của Kuon} \\[8pt]
             & {\caviar{\fontsize{14pt}{14pt}\selectfont Nhân vật: \emp{kuon2} \emp{kaigou2} \emp{suzuno2} \emp{naomi2} \emp{ryuuhou2} \emp{kson} \emp{delu} \emp{mike1} \emp{koneko} \emp{game}}} \\[8pt]
        \end{tabular}
    \end{adjustwidth}
\end{table}
%=========================================TRUYỆN CHÍNH========================================
\color{black}

\txtbx[2]{
    {\color{vol} \uline{Tóm tắt:}} Cuối tháng Tư, Kuon đột ngột sốt cao và phải nằm nghỉ ở tầng năm. Kaigou phát hiện anh vắng mặt ở bữa sáng rồi cuống cuồng chăm sóc; Suzuno nấu cháo, Naomi vẽ lời chúc, Ryuuhou nhận xử lý việc gấp còn Delutaya mang trà quỷ giới lên. Kson thay Kuon đến kho của mẹ anh, lần đầu tự tin đảm đương công việc bên ngoài mà không có anh bên cạnh. Đêm xuống, Mikeneko lặng lẽ để nước ấm và thuốc cạnh giường anh. Sáng hôm sau, Kuon không biết ai đã mang chúng tới. Anh hồi phục sau một ngày, đúng như Kaigou kể lại trong M2-04.
}

\stpt{-Phần 1-}
\stcap{Sát thủ thầm lặng}

\pov{Hikawa Kaigou}

\lang{Etsunan}

Cuối tháng Tư năm 2022, Kawachi đón một buổi sáng giao mùa với những thay đổi khó đoán, như thể thời tiết cũng đang lưỡng lự giữa mùa mưa vừa qua và mùa nắng sắp tới. Đêm trước đó, cả khu phố còn chìm trong mưa phùn lạnh lẽo, những hạt mưa nhẹ bết vào kính cửa sổ, nhưng chỉ trong vài giờ trước khi bình minh ló dạng, gió bỗng đổi chiều. Nó trở nên khô hơn, mang theo hơi ấm khó ngờ của một ngày sắp nắng, và luồng gió ấy làm những cánh cửa sổ của nhà Hikawa rung rinh từng nhịp nhỏ, như một tiếng gõ cửa thầm lặng gọi mọi người thức dậy.

Tôi thức dậy sớm đúng như thói quen hàng ngày, không có gì khác biệt so với mọi buổi sáng khác. Cơn buồn ngủ vẫn còn vương trên mi mắt, nhưng tôi đã đứng dậy, kéo rèm cửa một chút để ánh sáng đầu ngày lọt vào, thay bộ đồ ngủ bằng bộ đồ thường ngày ở nhà, dùng dây buộc gọn mái tóc vốn hay rủ xuống mặt. Sau khi chuẩn bị xong xuôi, tôi bước ra khỏi phòng, lên cầu thang tới phòng sinh hoạt chung, nơi cả nhà thường tập trung vào bữa sáng.

Cả căn nhà vẫn còn chìm trong sự yên tĩnh tuyệt đối, không có tiếng bước chân, không có tiếng mở tủ lạnh, thậm chí không có tiếng game nào phát ra từ phòng của các em. Điều đó không có gì lạ, Ryuuhou vốn hay ngủ muộn, dậy muộn, bình thường cũng phải gần sáu giờ mới lục đục lên bếp. 

Nhưng có một điều khiến tôi cảm thấy bất ổn, một chi tiết nhỏ mà tự nhiên tôi cảm nhận được ngay khi bước chân vào bếp: Aniki - anh Kuon - vẫn chưa xuất hiện.

Anh ấy luôn là người dậy sớm nhất trong nhà, kể cả những ngày cuối tuần không phải đi làm, kể cả những hôm có việc đêm trước mới ngủ muộn. Anh dậy trước tất cả chúng tôi, bước xuống bếp trước tiên, pha cho mình một tách trà xanh nóng, ngồi ở ghế sofa xem lịch công việc trong ngày, kiểm tra những cuộc hẹn, rồi mới chuẩn bị đồ đạc để đi làm.

Nhưng hôm nay, tách trà gỗ của anh vẫn đang úp trên giá treo ở bồn rửa, nguyên vẹn như chưa ai đụng đến. Chiếc cặp công việc màu đen mà anh hay mang đi làm vẫn nằm im lìm ở cạnh ghế sofa, đúng vị trí anh ấy vẫn để nó tối hôm trước, chứ không được đem ra cửa chuẩn bị gì cả. Thậm chí, đèn trong phòng anh ở tầng năm vẫn còn tắt, không có chút ánh sáng nào lọt ra khỏi khe cửa.

Tôi đứng ở chân cầu thang dẫn lên tầng năm, đôi chân dừng lại vài giây, lòng dâng lên một chút lo lắng khó nói. Tôi ngẩng đầu, hướng giọng nói lên trên tầng, vừa đủ để anh ấy nghe thấy nếu đã thức: ``Aniki? Anh dậy chưa?''

Không có một tiếng đáp nào vọng xuống. Cả căn nhà vẫn im lặng, chỉ còn tiếng gió thổi qua khe cửa, làm rung nhẹ bức tranh treo ở hành lang. Tôi gọi lần nữa, giọng lớn hơn một chút, để chắc chắn rằng nếu anh ấy đang trong phòng, anh sẽ nghe thấy: ``Anh Kuon? Dậy chưa ạ?''

Vẫn không có một chút động tĩnh nào. Ngay lúc đó, Ryuuhou bước ra từ phòng tắm ở cuối hành lang, mắt vẫn còn ngái ngủ, mí mắt vẫn sọng lại vì vừa mới thức dậy, tay vẫn đang cầm chiếc khăn lau mặt. Em nhìn tôi đứng ở chân cầu thang, nghiêng đầu hỏi: ``Chị gọi Onii à? Anh ấy chưa dậy sao?''

``Ừ. Hôm nay anh ấy chưa xuống bếp,'' tôi trả lời, mắt vẫn không rời khỏi cầu thang tối om.

Ryuuhou dụi mắt, nói một cách thản nhiên như thể đó là chuyện bình thường: ``Có khi anh ấy ngủ bù thôi, mấy hôm nay anh ấy có vẻ mệt mỏi lắm mà.''

Tôi quay đầu, nhìn chiếc đồng hồ treo trên tường phòng khách. Kim phút đã nhích qua số tám, đã quá sáu giờ rưỡi sáng, thời gian mà bình thường mà anh Kuon đã hoàn thành xong bữa sáng, đang chuẩn bị ra cửa để đi làm. Anh ấy hiếm khi ngủ muộn đến thế này, kể cả những ngày nghỉ lễ, anh ấy cũng không bao giờ nằm nhà đến giờ này mà không dậy. 

Tôi quay lại nhìn em gái, nói nhẹ nhàng: ``Em cứ xuống bếp ăn sáng trước đi. Chị lên xem anh ấy thế nào.''

Ryuuhou gật đầu đồng ý, nhưng em cũng không đi ngay. Em đứng yên ở đó, cùng tôi nhìn lên cầu thang tối om, ánh mắt cũng dâng lên một chút lo lắng như tôi. Tôi không nói gì thêm, đặt chân bước lên bậc thang đầu tiên, từng bước một leo lên tầng năm, nơi phòng của Kuon nằm.

\ct

Cánh cửa phòng của anh Kuon chỉ hé một khe nhỏ, như thể vừa có người bước ra hoặc còn đang giữ lại hơi ấm của căn phòng bên trong. Tôi đưa tay gõ nhẹ hai lần lên khung cửa, tiếng gõ nhỏ vừa đủ để nếu anh ấy đã tỉnh thì sẽ nghe thấy, trước khi đẩy cửa vào một cách nhẹ nhàng, cố gắng không làm phát ra tiếng động nào có thể làm phiền.

Bên trong phòng, tấm rèm cửa vẫn được kéo kín hoàn toàn, chặn hết mọi ánh sáng bên ngoài tràn vào, khiến không khí trong căn phòng trở nên ảm đạm và yên tĩnh đến lạ thường. Anh đang nằm nghiêng trên giường, một cánh tay đặt nặng lên trán, tư thế trông mệt mỏi vô cùng. Anh ấy vẫn mặc nguyên chiếc áo ngủ cũ từ tối hôm qua, còn tấm chăn thì bị đẩy xuống gần cuối giường, như thể cơ thể đang tỏa ra quá nhiều nhiệt khiến anh phải vật lộn thoát ra.

``Aniki? Anh nghe em nói không?'' tôi gọi nhỏ, bước lại gần cạnh giường, giọng nói vừa đủ để anh ấy nghe thấy.

Kuon mở mắt rất chậm, mí mắt nặng trĩu như có hàng tấn sức nặng đè lên, ánh mắt vẫn còn mơ màng chưa hoàn toàn tỉnh táo. Giọng nói của anh ấy vang lên khàn đặc, yếu ớt đến mức tôi phải cố gắng lắng nghe mới nhận ra được từ đầu tiên: ``Kaigou?''

Tôi đưa mu bàn tay lên đặt nhẹ lên trán của anh ấy, và ngay lập tức, một luồng nhiệt nóng rát truyền sang lòng bàn tay, khiến tôi lập tức rụt tay lại như bị bỏng. Cảm giác lo lắng trong lòng càng lúc càng tăng lên, tôi buột miệng: ``Anh nóng quá.''

Anh Kuon cố gắng chống tay lên nệm, gồng mình để ngồi dậy, cơ thể rung nhẹ vì sức lực không còn. Anh ấy vẫn không quên những công việc đang chờ đợi, giọng nói bật ra một cách khó khăn: ``Anh phải xuống dưới. Hôm nay có mấy việc cần xử lý, không thể nằm đây mãi được.''

Tôi nhanh chóng đặt tay lên vai của anh ấy, ấn nhẹ nhàng nhưng quyết liệt để anh ấy nằm lại xuống gối. Tôi lắc đầu, không để anh ấy có cơ hội chống đối: ``Không. Anh nằm yên đi, mọi việc đã có em và các em trong nhà lo liệu.''

``Chỉ cần xuống gọi điện thôi, xử lý vài cuộc họp quan trọng là được mà\dots'' anh ấy vẫn cố gắng phản đối, cố gắng ngồi dậy lần nữa.

``Anh đang sốt cao đấy, không thể lao vào công việc được.'' tôi nhấn mạnh từng chữ, hy vọng anh ấy sẽ nhận ra tình trạng nghiêm trọng của mình.

``Anh biết rồi\dots'' giọng Kuon yếu đi, có vẻ như tôi đã làm anh ấy chùn bước.

``Vậy thì đừng cãi nữa. Anh cần nghỉ ngơi,'' tôi nói, giọng nói vẫn giữ được sự cứng rắn nhưng ẩn chứa đầy sự lo lắng.

Kuon nhắm mắt lại, có lẽ định nói thêm điều gì đó để thuyết phục tôi, nhưng câu nói vừa đến miệng lại biến thành một tiếng ho khô, khè khè vang lên trong căn phòng yên tĩnh. Cơn ho khiến cả thân người anh ấy rung lên, làm tôi càng thêm hoảng sợ.

Tôi vội vàng quay sang ngăn tủ nhỏ ở đầu giường, nơi chúng tôi thường để các vật dụng y tế cơ bản, rút chiếc nhiệt kế thủy ngân ra. Tôi lau nhanh đầu nhiệt kế bằng một chiếc khăn giấy, rồi đưa cho anh ấy, đứng bên cạnh chờ đợi trong hồi hộp.

Ngay lúc đó, chiếc đồng hồ thông minh đặt trên bàn cạnh giường đột ngột sáng màn hình, Game khởi động và giọng nói điện tử bình thản của cậu ấy vang lên: ``Kaigou-user, tôi đã ghi nhận nhịp tim của Kuon-user cao hơn mức bình thường 30\%, nhiệt độ da tăng rõ rệt qua cảm biến. Hãy đo bằng nhiệt kế thủy ngân để có số liệu chính xác, đừng chỉ dựa vào cảm biến cổ tay.''

``Tôi đang làm đây!'' tôi đáp lại Game, mắt vẫn không rời khỏi anh Kuon đang cầm nhiệt kế trong tay.

Bàn tay của anh hơi run khi anh ấy cố gắng giữ nhiệt kế đúng vị trí, sức lực đã cạn kiệt sau cơn ho vừa rồi. Tôi nhanh chóng đưa tay đỡ lấy khuỷu tay của anh ấy, giữ cho nhiệt kế không bị xê dịch, cho đến khi thiết bị phát ra tiếng bíp quen thuộc báo hiệu đã đo xong.

Tôi cầm nhiệt kế lên, hướng về phía ánh sáng yếu ớt từ cửa để đọc con số, và tim tôi gần như ngừng đập khi thấy số liệu: ba mươi chín độ. Con số này cao hơn mức sốt thông thường, đủ để gây ra những biến chứng nguy hiểm nếu không được kiểm soát kịp thời.

Cảm giác tim tôi thắt chặt lại, một nỗi lo sợ lạ lùng tràn ngập trong lòng tôi. Tôi đã từng xử lý vô số tình huống khẩn cấp ở văn phòng, từ các cuộc họp đột xuất đến những vấn đề phát sinh trong công việc mà phải giải quyết trong vòng vài tiếng. Nhưng không tình huống nào khiến tôi hoảng loạn như lúc này, khi nhìn thấy người anh trai mình, người luôn mạnh mẽ và che chở cho cả nhà, đang nằm yếu ớt trên giường bệnh.

``Game, gọi phòng khám trực tuyến gần nhất ngay lập tức. Hỏi xem chúng ta cần làm gì để hạ sốt cho anh ấy!'' tôi ra lệnh, giọng nói không còn bình tĩnh như mọi khi.

``Đang kết nối với máy chủ của phòng khám\dots'' Game lập tức phản hồi, màn hình đồng hồ hiển thị thanh tải đang chạy.

Kuon hé mắt lên, cố gắng ngăn cản tôi: ``Không cần làm lớn chuyện đâu, chỉ là cơn sốt thông thường thôi, nghỉ ngơi một ngày là khỏi.''

Tôi đặt nhiệt kế xuống bàn cạnh giường với lực mạnh hơn dự định, tiếng động nhỏ vang lên trong căn phòng. Tôi nhìn thẳng vào mắt anh ấy, không chịu nhượng bộ: ``Anh đừng có nói câu đó khi đang sốt ba mươi chín độ được không? Đây không phải là chuyện nhỏ có thể bỏ qua được!''

Vị trợ lý ảo nhanh chóng kết nối thành công với phòng khám trực tuyến gần nhà, hình ảnh một nhân viên y tế mặc đồng phục hiện lên trên màn hình nhỏ. Tôi mô tả chi tiết tình trạng của Kuon: nhiệt độ đo được là 39 độ, có cơn ho khô, và dù mệt mỏi nhưng anh ấy vẫn hoàn toàn tỉnh táo.

Anh tôi vẫn có thể trả lời các câu hỏi của nhân viên y tế, không gặp khó khăn khi thở, cũng không có bất kỳ biểu hiện đau ngực hay chóng mặt nghiêm trọng nào. Nhân viên y tế sau khi lắng nghe đầy đủ thông tin đã hướng dẫn chúng tôi cần theo dõi sát tình trạng của anh ấy, cho Kuon nghỉ ngơi tuyệt đối, uống nước từng ngụm nhỏ thường xuyên và dùng thuốc hạ sốt đúng theo liều lượng ghi trên nhãn. Nếu xuất hiện bất kỳ dấu hiệu nặng hơn như khó thở, đau ngực hay nhiệt độ không hạ sau 24 giờ, chúng tôi phải gọi cấp cứu ngay lập tức.

Tôi ghi lại từng lời dặn của nhân viên y tế vào sổ tay nhỏ, không để sót bất kỳ chi tiết nào. Game cũng đồng bộ hóa toàn bộ hướng dẫn vào hệ thống, lưu lại trên màn hình để bất kỳ ai trong nhà cũng có thể xem xét. ``Tôi sẽ cài đặt nhắc nhở giờ đo nhiệt độ và uống nước, nhưng mọi quyết định y tế cuối cùng đều thuộc về nhân viên chuyên môn, tôi chỉ có thể hỗ trợ theo dõi.'' 

``Rõ rồi, cảm ơn ông,'' tôi đáp lại, gấp sổ tay lại.

Kuon nhắm mắt lại, giọng nói vẫn yếu ớt: ``Anh đã bảo là không cần gọi bác sĩ rồi, chỉ tốn thời gian và tiền bạc thôi.''

Tôi kéo tấm chăn lên ngang ngực cho anh ấy, chỉnh lại mép chăn để không có khoảng hở nào khiến anh ấy bị lạnh. ``Anh không được bảo gì hết, bây giờ việc duy nhất anh cần làm là nghỉ ngơi thật tốt để mau khỏe lại,'' tôi nói nhẹ nhàng, không còn cứng rắn như lúc nãy.

Tôi đặt một cốc nước lọc mới rót lên bàn cạnh giường, rồi bước đến cửa sổ, định mở hé tấm rèm để căn phòng có thêm một chút ánh sáng tự nhiên. Nhưng ngay khi ánh nắng đầu ngày tràn vào một phần căn phòng, Kuon đã nhăn mặt lại, rõ ràng là khó chịu với ánh sáng chói.

Tôi lập tức kéo tấm rèm lại một nửa, chỉ để lại một vệt sáng mỏng vừa đủ để căn phòng không quá tối, quay lại hỏi anh ấy: ``Sáng quá à? Em kéo lại rồi.''

``Ừ.'' anh ấy chỉ đáp được một từ ngắn ngủi.

``Được rồi, anh cứ nằm yên nghỉ ngơi. Em sẽ quay lại kiểm tra sau,'' tôi nói, rồi bước ra khỏi giường, tiến về phía cửa. Trước khi ra khỏi phòng, tôi lại ngoái đầu vào nhìn anh ấy một lần nữa, đảm bảo anh ấy đã nằm yên ổn.

Chỉ khi thấy anh Kuon vẫn nằm yên, nhắm mắt nghỉ ngơi, tôi mới thở phào nhẹ nhõm, bước chân nhẹ nhàng xuống cầu thang để báo lại tình hình với các em ở dưới.

\ct

Ở phòng khách, Suzuno đã dọn dẹp sạch sẽ, bày biện từng chiếc bát, đôi đũa và các món ăn sáng nóng hổi trên bàn, đủ cho tất cả mọi người trong nhà. Góc bàn vẫn còn một chiếc ghế trống, vị trí quen thuộc mà anh Kuon luôn ngồi mỗi sáng. Naomi, em út trong nhà, ngồi đối diện, mắt không rời khỏi chiếc ghế trống đó như đang chờ đợi một điều gì đó.

``Onii-chan chưa dậy ạ?'' em nhỏ giọng hỏi, đôi chân vẫn đung đưa nhẹ dưới chân ghế.

Tôi kéo ghế ngồi xuống, cật lực giữ giọng mình bình thản nhất có thể, nhưng trước khi kịp đưa tay cầm đôi đũa đã chuẩn bị sẵn, tôi đã phải nói ra sự thật khiến không khí bữa sáng lập tức thay đổi. ``Anh ấy bị sốt cao. Sáng nay mọi người nhớ nhẹ tay, đừng làm ồn để anh ấy yên nghỉ.''

Ngay lập tức, Ryuuhou - người vẫn đang dụi mắt, còn có vẻ ngái ngủ từ khi bước ra khỏi phòng tắm - hoàn toàn tỉnh táo, cơn buồn ngủ bay biến hết. Em đặt chiếc khăn lau mặt lên ghế, gương mặt dồn dập lo lắng: ``Sốt cao không chị? Bao nhiêu độ rồi ạ?''

``Ba mươi chín độ,'' tôi đáp, đến tôi cũng không tránh khỏi chút run rẩy dù đã cố gắng giữ bình tĩnh.

Naomi lập tức mở to mắt, đôi mắt tròn xoe ngẩn ngơ một lúc rồi vội vàng hỏi tiếp, giọng em nhỏ đi đáng thương: ``Onii-chan có đau lắm không ạ? Anh ấy có sao không?''

``Anh ấy chỉ mệt mỏi thôi, nhưng vẫn tỉnh táo, không có gì đáng lo ngại lắm đâu,'' tôi trả lời, cố gắng làm dịu đi nỗi hoảng loạn của em gái nhỏ, dù chính lòng mình vẫn còn đang cuộn lên từng cơn lo lắng.

Suzuno, người vừa cầm muôi chuẩn bị múc cơm cho mọi người, cũng đặt chiếc muôi xuống nhẹ nhàng, gương mặt em cũng thấm đẫm sự lo lắng. Em ngẩng lên nhìn tôi, nói một cách chắc nịch: ``Em nấu cháo cho anh ấy nhé? Cháo nóng dễ ăn, anh Kuon sẽ ăn được ít nhất vài thìa.''

Tôi gật đầu liền, lòng tràn đầy một chút nhẹ nhõm khi có em ấy lo liệu việc bếp núc. ``Nhớ nấu loại nhừ dễ ăn nhé, Suzu. Game đã liên hệ hỏi phòng khám trực tuyến rồi, bác sĩ bảo anh ấy cần uống nhiều nước và nghỉ ngơi tuyệt đối.''

``Vâng ạ, em hiểu rồi,'' Suzuno đáp lời, tay đã bắt đầu dọn dẹp các món ăn sáng để chuẩn bị nồi cháo mới.

Naomi cũng lập tức nhảy xuống khỏi ghế, bỏ lại bữa sáng chưa kịp ăn, chạy vội về phía hộp bút màu và tập giấy trắng em để ở góc phòng khách. Em vừa chạy vừa thốt lên vui vẻ như vừa nghĩ ra một ý tưởng tuyệt vời: ``Naomi sẽ vẽ tranh chúc Onii-chan mau khỏe! Em vẽ mặt trời và nhiều hoa thật đẹp cho anh ấy!''

Tôi định há miệng bảo em chạy chậm thôi, đừng vội để bị vấp chân, nhưng lời nhắc ấy lại mắc lại nơi cổ họng, không nói nên lời. Đôi mắt tôi bất giác quay sang nhìn lên cầu thang tối om dẫn lên tầng năm, nơi anh Kuon đang nằm nghỉ.`'

Đúng lúc đó, chiếc đồng hồ thông minh của Game đặt trên bàn rung nhẹ nhàng, màn hình sáng lên và giọng nói bình thản của cậu ấy vang lên, phá vỡ không khí im lặng ngắn ngủi: ``Kaigou-user, cô chưa ăn sáng gì cả. Hãy ăn ít nhất một chút đi, nếu không sức khỏe cũng sẽ bị ảnh hưởng.''

``Tôi sẽ ăn sau, đợi khi mọi việc ổn thỏa trước đã,'' tôi trả lời, mắt vẫn chưa rời khỏi cầu thang.

``Nếu cô ngất đi vì đói thì tôi sẽ phải theo dõi và nhắc nhở hai người cùng lúc, rất bất tiện. Vì vậy hãy ăn sáng trước đã,'' ông ta vẫn nhắc nhở một cách thẳng thắn. Phiền phức là vậy, nhưng ông ta nói nói không phải không có lý.

``Tôi không ngất đâu, đừng lo,'' tôi cười nhẹ, cuối cùng cũng quay mắt nhìn về chiếc đồng hồ trên bàn.

Suzuno bước tới từ phía bếp, mang theo một tách trà xanh nóng hổi, đặt nhẹ nhàng ngay trước mặt tôi. Em mỉm cười an ủi, nói nhỏ: ``Onee-sama uống chút nước ấm trước đi ạ. Uống xong cảm thấy sẽ khỏe hơn nhiều.''

Tôi đưa tay cầm chiếc tách trà lên, cảm nhận hơi ấm lan tỏa từ gốm sứ vào lòng bàn tay. Nhưng ngay cả khi hơi trà nóng bay lên, làm ấm cả bàn tay, tôi vẫn nhận ra rằng chính bàn tay của mình vẫn còn hơi run, không dứt được nỗi lo lắng về tình trạng của anh trai.

\ct

Tôi bước quay lại phòng của Kuon sau đó, bước chân vẫn nhẹ như lúc trước để không làm phiền bất cứ ai. Đặt lòng bàn tay lên trán anh, tôi lắng nghe cảm nhận nhiệt độ vẫn đang tỏa ra mạnh mẽ, rồi cầm lấy chiếc nhiệt kế đặt trên bàn để kiểm tra con số. Nhiệt độ vẫn chưa hề hạ đi chút nào, con số 39 độ vẫn hiện rõ trên bề mặt thủy ngân, khiến trái tim tôi lại thắt thêm một lần nữa.

Nghe tiếng bước chân của tôi, Kuon chậm rãi mở mắt, mí mắt vẫn nặng trĩu như thể mỗi lần mở mắt đều tốn rất nhiều sức lực.

``Em lại lên đây à?'' giọng anh vẫn khàn khàn, yếu ớt như trước.

``Em lên kiểm tra nhiệt độ cho anh.''

``Em mới kiểm tra lúc nãy thôi mà?''

``Đã mười lăm phút trôi qua rồi ạ.''

Kuon nghiêng người một chút, cố gắng nhìn chiếc đồng hồ treo trên tường, rồi thở ra một hơi nhẹ: ``\dots Ừ, đúng thật.''

Tôi kéo chiếc ghế gỗ nhỏ ở đầu giường lại gần cạnh giường, ngồi xuống để có thể nhìn rõ hơn tình trạng của anh. Rồi tôi bắt đầu hỏi những câu hỏi mà nhân viên y tế đã dặn phải theo dõi liên tục: ``Anh có thấy lạnh không?''

``Không.''

``Đau đầu nhiều không?''

``Có một chút.''

``Có thấy khó thở hay không?''

``Không.''

``Còn buồn nôn thì sao?''

``Cũng không nốt.''

Sau khi nhận đủ các câu trả lời, tôi thở phào nhẹ nhõm một chút; ít ra thì anh Kuon cũng không bị bệnh nặng. Tôi đẩy cốc nước lọc đã rót sẵn về phía anh: ``Được rồi. Anh uống ít nước đi.''

Kuon nhìn chiếc cốc nước trước mặt, rồi ngước mắt nhìn tôi, ánh mắt có chút bất lực trộn lẫn: ``Em có định hỏi anh thêm câu nào nữa không?''

Tôi khựng lại giữa lúc đang định sắp xếp mép chăn cho anh, những câu hỏi khác vẫn còn đang chen chúc trong đầu, nhưng lời anh nói đã làm tôi tỉnh ra. ``Nếu có thì sao?''

``Anh sẽ trả lời tất cả sau khi ngủ thêm một lúc nữa.'' Anh nói nhỏ, giọng đã có chút mệt mỏi.

Tôi tự nhận ra mình đang hỏi dồn anh, cứ như thể muốn lấp đầy mọi khoảng trống bằng câu hỏi để quên đi nỗi lo lắng đang cuộn trào trong lòng. Tôi hít một hơi thật sâu để bình tĩnh lại, đặt chiếc cốc nước chắc chắn vào tay anh, giọng nói cũng nhẹ nhàng hơn hẳn. ``Được rồi. Anh uống vài ngụm rồi ngủ tiếp nhé.''

Anh Kuon làm theo lời tôi, uống vài ngụm nước nhỏ rồi đặt cốc trở lại bàn cạnh giường. Sau đó, anh từ từ nằm xuống và quay mặt vào trong tường, quay lưng lại với tôi như một dấu hiệu rằng anh thật sự cần nghỉ ngơi.

Tôi ngồi yên trên ghế thêm một phút nữa, chỉ để lắng nghe nhịp thở của anh dần đều trở lại, đảm bảo anh đã thực sự chìm vào giấc ngủ. Ngay lúc đó, chiếc đồng hồ thông minh của Game đặt trên bàn bật sáng nhẹ, giọng nói của cậu ấy nhỏ đến mức chỉ có thể nghe thấy nếu ngồi đủ gần: ``Kuon-user đã ngủ rồi. Cô cũng nên xuống dưới ăn sáng đi, tôi đã nhắc nhiều lần mà.''

Tôi đứng dậy từ chiếc ghế, chỉnh lại quần áo một chút rồi đáp nhỏ: ``Biết rồi.''

Tôi bước ra khỏi phòng, khép cánh cửa lại một cách nhẹ nhàng nhất có thể, để không có một tiếng động nào có thể làm phiền giấc ngủ của anh.

\ct

Sau khi khép cửa phòng Kuon lại, tôi bước chậm rãi xuống cầu thang, bước chân vẫn cố gắng nhẹ nhàng nhất để không làm phiền giấc ngủ của anh. Khi vừa đặt chân đến bếp, tôi thấy Suzuno đang mở cửa tủ gỗ, lấy ra một túi gạo thơm và một chiếc nồi nhỏ bằng inox, đặt nhẹ lên mặt bếp đá granit màu xám. 

Naomi đứng sát bên cạnh chị, hai tay ôm chặt một tờ giấy trắng lớn cùng chiếc hộp bút màu nhựa em thường dùng, mắt lấp lánh theo dõi từng động tác của mọi người. Ryuuhou kéo một chiếc ghế gỗ thấp từ góc phòng, đẩy nó đến gần quầy bếp để em út có thể đứng lên nhìn rõ mọi thứ đang diễn ra trong bếp.

``Naomi vẽ xong chưa?'' Suzuno nhẹ nhàng hỏi, vừa lấy chiếc bình đựng gạo ra.

``Naomi sắp xong rồi ạ,'' em nhỏ đáp, đầu vẫn cúi xuống vẽ nhẹ trên tờ giấy.

Tôi bước lại gần để nhìn bức tranh của em: trong đó, Kuon nằm trên chiếc giường quen thuộc ở tầng năm, bên cạnh có chiếc cốc nước gỗ mà anh vẫn dùng, phía trên là một mặt trời màu vàng tươi vẽ lớn, tỏa sáng khắp góc tranh. Phía trên đầu anh trong bức vẽ, em đã viết bằng bút chì màu đậm dòng chữ nhỏ: ``Onii-chan mau khỏe nhé.''

Vừa nói xong, Naomi lại cầm bút màu đỏ, vẽ thêm một trái tim nhỏ xinh ở góc dưới cùng của tờ giấy, như muốn gửi thêm tình yêu thương của mình cho anh.

Ryuuhou dựa lưng vào quầy bếp, tay khoanh trước ngực, giọng nói trầm ổn: ``Em nhận xử lý mấy cuộc gọi công việc cho Onii. Đã nhắn với A-chan là hôm nay anh ấy nghỉ, không thể tham gia các cuộc họp được rồi.''

Tôi vừa bước xuống cầu thang xong, vừa tình cờ nghe được câu nói của em. Tôi tiến lại gần, hỏi nhẹ: ``Còn mấy việc gấp cần xử lý không?''

``Em đã chuyển phần việc khẩn cấp cho người phụ trách khác nhận rồi ạ. Việc nào có thể trì hoãn thì để mai xử lý tiếp,'' Ryuuhou đáp, giọng vẫn bình thản.

``Tốt lắm. Đừng tự mình ôm hết mọi việc lên người quá,'' tôi nói, lo lắng em sẽ quá sức.

Ryuuhou nhún vai, cười nhẹ: ``Em đâu có ôm hết đâu. Em chỉ chia việc cho mọi người cùng lo thôi mà.''

Đúng lúc đó, từ phòng khách vang lên tiếng bíp nhỏ của chiếc đồng hồ thông minh của Game, cậu ấy đã kích hoạt chế độ nhắc nhở. ``Tôi đã chuyển trạng thái lịch của Kuon-user sang nghỉ bệnh. Tất cả các cuộc họp sáng nay đều đã được dời sang ngày khác hoặc hủy bỏ,'' giọng nói điện tử của Game vang lên rõ ràng.

Tôi liếc nhìn chiếc đồng hồ đặt trên bàn trà, rồi nói: ``Cảm ơn ông nha Game.''

``Đừng cảm ơn. Tôi làm vậy để cậu ấy không cố gắng tham gia cuộc họp bằng giọng của người sắp mất tiếng đâu,'' ông ta đáp lại vẫn giữ giọng thẳng thắn quen thuộc.

Ở bếp, Suzuno đã rửa sạch gạo dưới vòi nước, đong nước đúng tỷ lệ để nấu cháo rồi đặt chiếc nồi lên bếp, bật lửa nhỏ liu riu. Naomi đứng trên chiếc ghế thấp Ryuuhou vừa mang ra, mắt tròn xoe quan sát từng động tác, như muốn học cách nấu cháo ngon như em ấy. Ryuuhou mở ngăn đựng gia vị, lấy ra hộp muối ăn và một củ gừng tươi mới mua từ chợ sáng hôm trước, đặt lên thớt.

Trước khi dùng bất kỳ thứ gì, Suzuno đều lật từng nhãn trên các lọ gia vị, kiểm tra kỹ lưỡng để đảm bảo không lấy nhầm loại. Tôi để ý đến động tác cẩn thận ấy của em, mỉm cười: ``Em cẩn thận thật.''

``Em muốn chắc chắn mọi thứ đều đúng loại, không có gì sai sót ạ,'' Suzuno đáp, tay vẫn tiếp tục công việc của mình.

Sau đó, em đặt củ gừng lên thớt, dùng dao sắc cắt thành những lát mỏng đều tăm tắp. Mùi thơm dịu đặc trưng của gừng tươi lan tỏa khắp căn bếp, làm không khí ấm áp hơn bao giờ hết.

\ct

Trong khi nồi cháo trên bếp sôi liu riu, hơi nước bay lên nhẹ nhàng, Suzuno chuẩn bị thêm một ít hành lá tươi, rửa sạch rồi thái nhỏ để vào một chiếc đĩa riêng, chuẩn bị cho khi ăn cháo. Naomi đặt bức tranh vừa vẽ xong vào một chiếc bìa cứng dày, để tờ giấy không bị cong hay nhàu trong khi mang lên cho Kuon. 

Ryuuhou mang một tập hồ sơ dày từ phòng làm việc xuống, đặt lên bàn trà trong phòng khách: ``Đây là mấy lịch họp hôm nay. Em đã đánh dấu việc nào cần người thay thế xử lý, để mọi người dễ theo dõi.''

Tôi tiến lại, nhanh chóng lật từng trang hồ sơ để xem xét, bất ngờ hỏi: ``Em còn ghi cả số điện thoại của người phụ trách dự phòng cho mỗi việc à? Chu đáo thật.''

``Em đã gọi điện xác nhận với họ trước rồi, mọi người đều đồng ý nhận việc thay thế cho anh Kuon rồi ạ,'' em đáp, giọng có chút tự hào nhỏ.

Tôi ngẩng đầu lên, nhìn em gái với ánh mắt ngạc nhiên: ``Em làm hết những chuyện này trong lúc mọi người vừa mới ăn sáng á? Thời gian ngắn thế mà em xử lý được hết?''

``Em đã bắt đầu chuẩn bị từ trước khi chị gọi xuống bếp báo tin anh ấy bị ốm rồi ạ,'' em nói, như đã dự đoán trước tình hình.

Ánh mắt tôi tự nhiên dịu đi khi nhìn em, lòng tràn đầy sự biết ơn: ``Cảm ơn em đã lo liệu mọi thứ chu đáo thế.''

Ryuuhou gãi má, ngại ngùng cười: ``Em chỉ không muốn anh ấy thức dậy rồi lại phải lo lắng về công việc, mà quên mất việc nghỉ ngơi để mau khỏe thôi.''

Tôi không trêu chọc em như mọi khi, chỉ im lặng nhìn em. Tôi gấp tập hồ sơ lại cẩn thận, đặt nó lên góc bàn trà để sau này có thể xem lại. Lúc này, ở bếp, Suzuno múc một ít cháo ra bát nhỏ, kiểm tra độ đặc của nước cháo. Hạt gạo đã nở mềm hoàn toàn, nước cháo không quá loãng mà cũng không quá sánh, đúng độ ngon nhất.

Em nếm một thìa nhỏ để thử vị, rồi thêm chút muối vào nồi cháo để vừa ăn. Naomi nhanh tay đưa cho chị một chiếc thìa sạch khác, để em tôi không cần dùng lại chiếc thìa vừa nếm. 

Suzuno cười rạng rỡ, nói với em nhỏ: ``Cảm ơn Naomi-chan, em chu đáo thật.''

Tôi bước lại gần bếp, ngửi thử mùi thơm của nồi cháo đang sôi, mùi gạo nướng cùng gừng tươi lan tỏa rất dễ chịu: ``Thơm đấy, chắc chắn anh ấy sẽ ăn được.''

``Em sẽ mang lên một phần nhỏ trước. Nếu anh ấy ăn được thì em mới mang thêm phần khác lên sau ạ,'' Suzuno nói, vẫn cẩn thận với tình trạng của anh Kuon.

Game ngay lập tức xác nhận lại hướng dẫn của phòng khám: ``Phòng khám không yêu cầu ép người bệnh phải ăn, chỉ cần đảm bảo cậu ấy uống đủ nước và nghỉ ngơi tuyệt đối là được.''

Suzuno chia cháo ra thành một phần nhỏ, đậy nắp khay nhựa giữ nhiệt để mang lên. Naomi cầm bức tranh đã được bọc bìa cứng, đi theo sau chị. Tôi là người bước cuối cùng, tay mang theo chiếc nhiệt kế thủy ngân để có thể kiểm tra nhiệt độ cho anh Kuon ngay khi lên phòng.

\ct

Lúc Suzuno gõ cửa phòng của anh Kuon, anh ấy đã chìm vào giấc ngủ từ lâu. Tôi là người bước tới trước, nhẹ nhàng hé cánh cửa chỉ đủ để nhìn lướt qua bên trong, tránh làm phát ra tiếng động nào có thể đánh thức anh. Anh vẫn nằm nghiêng như lúc tôi rời đi, nhịp thở đều đặn, tưởng chừng như mọi cơn sốt đều đang dịu đi trong giấc ngủ.

Em tôi bước vào sau, đặt nhẹ chiếc khay giữ nhiệt lên chiếc bàn thấp cạnh giường, nơi vẫn còn đặt cốc nước lọc từ sáng. Naomi đi theo sát sau, nhỏ chân tiến lại gần nhưng dừng lại cách giường một khoảng an toàn, như sợ chỉ cần bước thêm một bước cũng có thể làm phiền giấc ngủ của anh trai.

``\hl{Onii-chan đang ngủ ạ?}'' em thì thầm, giọng nhỏ đến mức gần như tan vào không khí ảm đạm của căn phòng.

``\hl{Ừ. Mình để anh ấy ngủ thêm một chút,}'' Suzuno cũng đáp lại bằng một giọng nhỏ không kém, đặt ngón tay lên môi ra hiệu giữ yên lặng.

Naomi gật đầu hiểu ý, nhẹ nhàng đặt bức tranh em đã vẽ cả buổi sáng dựa vào chiếc cốc nước trên bàn, để anh Kuon có thể nhìn thấy ngay khi tỉnh dậy. 

Tôi bước lại phía bàn, cầm lấy chiếc nhiệt kế thủy ngân để kiểm tra lại nhiệt độ cho anh, trái tim tôi thắt lại khi con số hiện ra vẫn không thay đổi nhiều: gần ba mươi chín độ, cơn sốt vẫn chưa chịu lui đi chút nào. Tôi nhanh chóng ghi lại con số này vào bảng theo dõi mà Game đã mở sẵn trên màn hình đồng hồ, để mọi người trong nhà đều có thể nắm được tình hình của anh.

Ngay lúc đó, anh Kuon khẽ cựa mình, rồi từ từ mở mắt, mí mắt vẫn nặng trĩu như chưa thể thoát khỏi cơn mệt mỏi. ``Suzuno\dots?'' anh gọi khẽ, nhận ra bóng dáng của Suzuno bên cạnh giường.

``Em mang cháo lên cho anh. Anh có muốn thử vài thìa không ạ?'' em đáp lại nhỏ nhẹ, chỉ để cho anh Kuon nghe thấy.

Anh nhìn vào bát cháo nóng hổi trên khay, ánh mắt vẫn còn mơ màng, rồi mới đáp: ``Một chút thôi.''

Tôi nhanh chóng bước tới, kê một chiếc gối dày phía sau lưng anh để anh có thể ngồi dậy một cách thoải mái, không cần phải gồng mình lên. Suzuno cầm lấy bát cháo, dùng thìa múc chỉ một lượng nhỏ vừa đủ, từ từ đưa đến cho anh để anh không bị sặc. Kuon ăn rất chậm, mỗi thìa như phải tốn rất nhiều sức lực, khiến lòng tôi càng thêm xót xa.

Naomi vẫn đứng nép gần cửa, ôm chặt hai tay trước ngực như đang hồi hộp chờ đợi, rồi em nhỏ giọng lên tiếng: ``Em vẽ tranh cho anh đó ạ.''

Kuon ngẩng đầu nhìn về phía bức tranh trên bàn, khóe miệng yếu ớt nhếch lên thành một nụ cười mệt mỏi, nhưng rõ ràng là vui vẻ. ``Cảm ơn nhé Naomi. Mặt trời trong tranh sáng thật.''

``Naomi muốn anh mau khỏe,'' em đáp lại nhanh chóng, đôi mắt tròn xoe lấp lánh hy vọng. ``Khi nào anh khỏe rồi, Naomi muốn anh chơi với em nhiều hơn ạ!''

``Ừm.''

Tôi lại khẽ chạm lòng bàn tay lên trán của anh lần nữa, cảm nhận luồng nhiệt vẫn tỏa ra mạnh mẽ, thì anh Kuon khẽ liếc sang nhìn tôi. ``Em lại kiểm tra à?''

``Em đang xác nhận nhiệt độ thôi,'' tôi đáp lại, cố gắng giữ giọng mình bình thản nhất.

``Nghe giống nhau mà,'' anh cười khẽ, vẫn là giọng nói yếu ớt.

Suzuno đứng bên cạnh nhìn tôi nhưng không nói gì, chỉ lặng lẽ theo dõi tình hình của anh trai. Tôi đặt nhiệt kế trở lại trên bàn, rồi nhẹ nhàng chỉnh lại mép chăn cho anh, đảm bảo không có khoảng hở nào khiến anh bị lạnh. 

``Ăn xong thì ngủ tiếp nhé. Mọi việc bên ngoài đều có người xử lý rồi, anh không cần phải lo lắng gì cả.''

Kuon gật đầu, không hỏi thêm gì nữa, như thể đã đặt  hoàn toàn tin tưởng vào chúng tôi. Anh ăn thêm hai thìa cháo nhỏ nữa, uống vài ngụm nước ấm rồi từ từ nằm xuống, chìm lại vào giấc ngủ. Naomi trước khi rời đi cũng vẫy tay thật khẽ với anh trai, như một lời chúc ngủ ngon.

\ct

Khi mặt trời dần leo lên cao, đến gần giữa trưa, cơn sốt của Kuon vẫn chưa hề hạ nhiệt.

Tôi đã lẻn quay lại phòng anh ấy tổng cộng ba lần chỉ trong vòng mấy tiếng ngắn ngủi.

Lần đầu tiên, tôi mang theo chiếc nhiệt kế thủy ngân cũ để kiểm tra lại con số trên đó, mong mỏi thấy nhiệt độ đã giảm đi dù chỉ nửa độ.

Lần thứ hai, tôi mang thêm một cốc nước lọc mới rót, sợ anh thức dậy mà khô cổ, không có gì để uống.

Lần thứ ba, tôi chỉ ngồi ở mép giường, nhẹ nhàng hỏi anh có thấy rét run hay không, vì cơn sốt cao thường khiến người bệnh lạnh lùng từ trong xương.

Game bật giọng bình thản từ chiếc đồng hồ thông minh đặt trên bàn khách, đọc lại một lần nữa toàn bộ lời hướng dẫn của nhân viên y tế, nhắc nhở cả nhà phải ghi lại nhiệt độ của Kuon đúng theo khoảng thời gian đã được dặn, không nên kiểm tra quá thường xuyên cũng không được để quá lâu.

Tôi gật đầu nghe theo, nhưng cứ mỗi lần chân tôi đi ngang qua cánh cửa phòng tầng năm, trái tim tôi lại như bị kéo lại, không tự nhiên mà dừng chân ở cửa, muốn hé nhìn vào bên trong một chút nữa.

Rồi, Ryuuhou mang chiếc điện thoại của mình xuống phòng khách, màn hình vẫn còn sáng lên dòng tin nhắn mới. ``Kson-nee vừa gọi rồi chị ạ. Chị ấy nói sẽ đến kho thay Onii xử lý công việc hôm nay mà anh ấy đã hứa với mẹ hôm kia.''

Tôi ngẩng đầu khỏi chiếc ly nước đang cầm trên tay, bất ngờ hỏi: ``Kson tự mình đến kho à? Không có ai đi cùng à?''

``Dạ. Mẹ đã gặp em ấy sáng nay, nói rõ hết những công việc cần làm, những số liệu cần kiểm tra rồi ạ.''

Tôi thở ra một hơi nhẹ nhõm, rồi gật đầu với em. ``Tốt lắm. Nhắn lại với chị ấy rằng đừng gọi điện cho anh Kuon trừ khi có chuyện thật sự cấp bách, không thể giải quyết được mà thôi.''

Game ngay lập tức góp ý thêm, giọng điện tử vẫn đều đều: ``Kho đã có nhân viên phụ trách tại chỗ rồi, cô ấy chỉ cần tiếp quản phần lịch giao hàng và kiểm tra tồn kho vật tư thôi, không có gì quá phức tạp cần hỏi Kuon-user cả.''

Ryuuhou nhanh chóng gõ trả lời tin nhắn trên điện thoại, rồi ngẩng đầu nhìn lên bậc cầu thang tối om dẫn lên tầng năm, ánh mắt có chút lo lắng nhìn tôi.

``Chị Kaigou này, chị đã không nghỉ ngơi gì cả từ sáng nay rồi. Chị có thể ngồi xuống nghỉ một chút không ạ?''

``Chị ổn mà, không mệt đâu.'' Tôi đáp lại một cách tự nhiên.

``Chị nói câu đó y hệt như anh Kuon luôn ấy. Lúc nào cũng nói mình ổn, dù có mệt mỏi đến đâu cũng không chịu thừa nhận,'' em cười khẽ, nhưng nụ cười đó có chút lo lắng.

Tôi nhìn em, không biết nói gì để đáp lại.

Con bé lập tức quay đi, tránh cho tôi phải bối rối, vừa bước về phía bếp vừa nói: ``Em lên bếp xem chị Suzu cần giúp gì thêm đây.''

Tôi há miệng định gọi em lại, muốn nói rằng tôi không có ý gì, chỉ là tôi thật sự không mệt, nhưng rồi lại thôi, để em đi.

Tôi kéo chiếc ghế gỗ nhỏ đặt bên ngoài cửa phòng Kuon, ngồi xuống đó, chạm vào bề mặt gỗ ấm áp dưới lòng bàn tay.

Lần này, tôi tự nhủ phải bình tĩnh, cố đợi đúng khoảng thời gian Game đã nhắc nhở, không thể cứ chạy lên chạy xuống làm phiền giấc ngủ của anh nữa.

\ct

\lang{Nhật}

Đến đầu giờ chiều, khi mặt trời đã leo cao hơn và ánh nắng dịu đi một chút, Delutaya bước lên cầu thang, hai tay nâng niu một chiếc bình nhỏ được bọc kín bằng khăn bông dày. Khăn bọc giữ lại hơi ấm của thứ bên trong, để lại chút hương thơm thoang thoảng bay theo bước chân cô.

``Em pha trà theo công thức cũ của nhà em,'' cô nói nhỏ, đặt chiếc bình nhẹ nhàng xuống bàn trong phòng khách trước khi tôi có thời gian hỏi. ``Trong đó có gừng tươi, lá bạc hà mới hái và một chút mật ong. Nếu anh không cảm thấy muốn uống, hay cổ họng anh khó chịu để nuốt, mình để nguội hoàn toàn rồi cũng được.''

Tôi dán mắt vào chiếc bình bọc khăn, lòng vẫn không khỏi băn khoăn dù biết cô nàng chưa bao giờ làm gì không chín chắn.

``Em có hỏi qua phòng khám xem những nguyên liệu này có thể dùng cùng lúc với thuốc anh ấy đang uống không chưa?'' tôi hỏi, giọng vẫn giữ được sự thận trọng suốt ngày nay.

Delutaya lập tức gật đầu, không chần chừ: ``Em đã hỏi Suzuno-san trước, em ấy kiểm tra lại với hướng dẫn bác sĩ gửi rồi. Đây chỉ là trà pha loãng thôi, không có thêm thuốc hay bất kỳ thành phần lạ nào có thể gây ảnh hưởng ạ.''

Ngay lập tức, chiếc đồng hồ thông minh của Game trên bàn cũng phát ra tiếng xác nhận, giọng điện tử bình thản như mọi khi: ``Tôi đã kiểm tra danh sách nguyên liệu, không có thứ gì xung đột với thuốc Kuon-user đang dùng. Các thành phần đều an toàn.''

Tôi thở ra một hơi dài mà chính mình cũng không nhận ra đã nín nhịn từ bao giờ, vai vốn căng cứng từ sáng nay mới thả lỏng được một chút.

``Được rồi. Nhưng chỉ để anh ấy uống nếu anh ấy thực sự muốn thôi, đừng ép,'' tôi nhắc lại lần nữa, quen với việc anh trai chúng tôi luôn cố gắng đáp lời mọi người dù bản thân không thoải mái.

Delutaya gật đầu hiểu ý, cô cũng quen với tính cách của anh trai chúng tôi.

Ngay khi chúng tôi vừa nói nhỏ mấy câu, tiếng trở mình nhẹ vang lên từ trong phòng tầng năm. Anh Kuon đã tỉnh dậy khi nghe thấy tiếng nói bên ngoài cửa phòng.

Anh không mở mắt ngay, có lẽ mí mắt vẫn nặng trĩu vì cơn sốt, chỉ hơi thở có chút thay đổi để lộ ra mình đã thức.

Tôi bước vào phòng trước, bước chân vẫn nhẹ nhàng như mọi lần để không làm anh giật mình.

``Aniki, em mang trà lên. Anh có muốn thử một ngụm để làm ấm cổ họng không?'' tôi hỏi nhỏ, tiến lại cạnh bàn cạnh giường.

Kuon từ từ mở mắt, ánh mắt còn mơ màng, nhìn chiếc bình trà đặt trên bàn rồi mới đáp lời, giọng vẫn còn khàn nhưng có vẻ đỡ hơn lúc sáng một chút: ``Một ngụm thôi.''

Delutaya bước vào sau tôi, cẩn thận mở nắp chiếc bình và rót trà ra chiếc cốc gỗ nhỏ Kuon vẫn dùng hàng ngày. Hơi nóng bay lên, mang theo hương bạc hà dịu nhẹ hòa quyện cùng mùi gừng ấm áp, lan tỏa khắp căn phòng vốn còn hơi ảm đạm.

Kuon nâng chiếc cốc lên, nhấp nhẹ một ngụm nhỏ rồi từ từ nuốt, sau đó gật đầu với chúng tôi. ``Vị dễ chịu lắm, cổ họng đỡ khô hơn nhiều rồi.''

Delutaya mỉm cười rạng rỡ, giọng đầy vui mừng: ``Em mừng vì anh thích.''

Lúc đó, tôi chợt nhớ ra chưa kiểm tra nhiệt độ cho anh, liền với tay lấy chiếc nhiệt kế thủy ngân đặt trên bàn. Tôi đưa mắt nhìn con số hiện trên que thủy ngân, và trái tim tôi bất giác rung lên. Con số đã giảm một chút so với lần kiểm tra cách đây nửa tiếng, không còn đứng yên ở mốc 39 độ nữa.

Tôi cố gắng không để lộ vẻ nhẹ nhõm ra mặt, sợ anh thấy mà lo lắng rằng mình đã quá hồi hộp, nhưng chính bàn tay đang cầm chiếc nhiệt kế lại bộc lộ tất cả - những ngón tay tôi siết chặt vào que thủy ngân lạnh, như muốn ghi nhớ con số hy vọng đó vào lòng.

\ct

Ryuuhou dành cả buổi chiều để hoàn thành những công việc mà Kuon đã phân công từ trước khi bị ốm. Em ngồi ở bàn trà, chiếc điện thoại được đặt ngay bên cạnh để gọi lần lượt cho hai người phụ trách chính, thông báo về việc dời lịch các cuộc họp và gửi qua email toàn bộ tài liệu liên quan cần thiết để họ có thể nắm bắt. 

Sau khi hoàn thành tất cả, tôi lần lượt kiểm tra lại từng việc một, đọc từng ghi chú em để lại và xác nhận các cuộc gọi đã được thực hiện. Không một mục nào trong danh sách bị bỏ sót, mọi thứ được sắp xếp gọn gàng, rõ ràng.

``Em làm tốt lắm,'' tôi nói, gật đầu hài lòng với sự chu đáo của em gái.

Ryuuhou chỉ nhún vai một cách tự nhiên, như thể đó chỉ là những việc nhỏ nhặt đáng làm: ``Em chỉ làm theo đúng danh sách đã được lập sẵn thôi.''

``Nhưng danh sách đó cũng phải có người sắp xếp, lọc lẻ các việc quan trọng ra chứ,'' tôi nhấn mạnh.

``Game đã giúp em lọc những việc cần ưu tiên gấp mà,'' em giải thích.

Ngay lập tức, chiếc đồng hồ thông minh trên bàn bật sáng, Game lên tiếng xác nhận: ``Tôi chỉ sắp xếp thứ tự ưu tiên thôi. Chính Ryuuhou-user đã tự mình thực hiện tất cả các cuộc gọi và xác nhận lại với từng người.''

Tôi lại gật đầu một lần nữa, lòng tràn đầy sự tin tưởng vào em gái. ``Vậy thì cả hai đều đã làm rất tốt.''

Ryuuhou nhìn tôi, rồi mỉm cười rạng rỡ, nụ cười ấy thổi bay đi chút lo lắng vẫn còn vương vấn trong lòng tôi từ sáng nay. Ở góc phòng khách, Naomi ngồi bệt xuống sàn, cây bút màu xanh lá cây di chuyển đều đặn trên tờ giấy vẽ. Em đang tô thêm một chiếc chăn màu xanh ấm áp quanh hình ảnh Kuon trong bức tranh của mình, như muốn gửi thêm hơi ấm cho anh trai đang nằm nghỉ.

Ở trong bếp, Suzuno đang kiểm tra lại nồi cháo còn lại trên bếp, em múc phần cháo chưa ăn hết vào những chiếc hộp giữ nhiệt để cất trong tủ lạnh, đảm bảo nếu anh Kuon đói vào buổi tối vẫn có thể hâm nóng lại dễ dàng. 

Delutaya vừa rửa sạch chiếc bình trà mà cô mang đến chiều nay, lau khô rồi đặt nhẹ nhàng lên chiếc khay gỗ ở góc bếp để chuẩn bị cho lần sử dụng tiếp theo.

Tôi lại bất giác ngước mắt nhìn lên bậc cầu thang tối om dẫn lên tầng năm, nơi cánh cửa phòng anh Kuon vẫn đang đóng kín. Ngay lúc đó, giọng nói đều đều của Game vang lên, như đọc được suy nghĩ trong đầu tôi: ``Cô đã kiểm tra nhiệt độ cho Kuon-user cách đây mười hai phút rồi.''

Tôi khoanh tay trước ngực, thừa nhận sự thật mà không thể chối cãi. ``Biết chứ.''

``Vậy thì cô phải đợi thêm ít nhât ba phút nữa trước khi kiểm tra lại lần tiếp theo,'' Game nhắc nhở đúng theo quy trình mà nhân viên y tế đã hướng dẫn.

Tôi liếc nhìn chiếc đồng hồ thông minh trên bàn, nhìn vào bộ đếm thời gian đang chạy ngược. ``Được.''

Và thật kỳ diệu, tôi đã thực sự ngồi yên lặng trên ghế suốt ba phút đó, không đứng dậy, không lẻn lên cầu thang, không cố gắng hé cửa để nhìn vào bên trong phòng như bao lần trước. Ryuuhou đã nhận ra sự nỗ lực giữ bình tĩnh của tôi, nhưng em chỉ im lặng quan sát chứ không hề trêu chọc hay nhắc đến điều đó.

\ct

Buổi chiều hôm đó, điện thoại của tôi vang lên, và khi tôi nhấc máy, giọng của Kson từ kho vọng ra rõ ràng qua chế độ loa ngoài mà tôi vừa bật. Tôi vừa quay lại ghế sofa sau khi lướt qua kiểm tra cửa phòng Kuon một lần nữa, nên mọi người trong phòng đều có thể nghe được cuộc trò chuyện.

``Kaigou-paisen, chị nghe rõ không?'' giọng cô ấy vang lên, có chút hơi thở nhanh như vừa chạy xong một công đoạn nào đó.

``Nghe rõ lắm. Mọi thứ ở đó vẫn chứ?'' tôi hỏi, tay vô thức siết chặt vào cạnh điện thoại, vẫn còn lo lắng liệu cô ấy có gặp trở ngại gì trong công việc hay không.

``Hoàn toàn ổn ạ. Em vừa kiểm tra xong toàn bộ lịch giao hàng cho ngày hôm nay, đã ký nhận lô vật tư mới đến từ nhà cung cấp, và cũng đã dời tất cả các cuộc hẹn của Onii sang ngày mai rồi ạ.''

Tôi thở nhẹ một hơi, rồi hỏi tiếp: ``Có vấn đề gì phức tạp mà vẫn cần anh ấy quyết định không? Bất cứ việc gì nếu vượt quá phạm vi hay cần sự phê duyệt đặc biệt thì cứ nói, bọn chị sẽ tìm cách xử lý.''

``Chỉ có một đơn hàng cần đổi thứ tự giao thôi ạ. Em đã hỏi người phụ trách tuyến vận tải rồi, họ xác nhận là có thể sắp xếp lại lộ trình để xử lý được, không có gì trở ngại cả.''

Tôi tự nhiên quay đầu nhìn sang Game, chiếc đồng hồ thông minh của anh Kuon đang nằm trên bàn trà, màn hình vẫn sáng nhè nhẹ. Ông ta cũng lập tức hiểu ý, nhanh chóng truy cập vào hệ thống lịch chung của kho để kiểm tra thông tin như em ấy vừa nói, rồi gật đầu với tôi để xác nhận mọi thứ đều khớp.

Tôi quay lại với điện thoại, đáp lại Kson: ``Vậy em tự quyết định như vậy là đúng rồi. Không có gì phải lo lắng nữa.''

Đầu dây bên kia im lặng một nhịp ngắn, như thể em ấy đang nghỉ ngơi sau buổi sáng bận rộn, rồi giọng cô ấy lại vang lên, có chút nhẹ nhõm: ``Em chỉ muốn báo cho mọi người ở nhà biết là em xử lý được tất cả mọi việc thôi. Không có gì bị sót lại hay bị trì hoãn đâu.''

``Chị biết em làm được.'' tôi nói thật lòng, lòng tự nhiên thấy ấm áp khi biết Kson đã trưởng thành đến thế.

``Còn Kuon-paisen thì sao? Tình hình hiện tại thế nào rồi ạ?''

Tôi vô thức ngước mắt nhìn về phía cánh cửa phòng tầng năm vẫn đang đóng kín, nơi anh Kuon vẫn đang chìm trong giấc ngủ nghỉ ngơi. Giọng tôi cũng nhẹ đi khi trả lời: ``Anh ấy vẫn đang nghỉ ngơi. Nhiệt độ của anh ấy bắt đầu hạ rồi, không còn giữ ở mức cao như lúc sáng nữa.''

Tiếng thở phào nhẹ nhõm vang lên ở đầu dây bên kia. ``Tốt quá rồi. Chị nhớ đừng để anh ấy nghe máy bất cứ lúc nào đâu, dù có chuyện gì đi nữa. Anh ấy cần nghỉ ngơi hoàn toàn, không được phép lao vào công việc khi vẫn còn ốm.''

``Chị đang cố gắng giữ anh ấy ở lại giường mà.'' tôi đáp lại, không khỏi bật cười nhỏ vì sự cẩn thận của Kson.

``Chỉ đang cố thôi à?'' Kson đáp lại bằng câu hỏi y hệt như những gì anh Kuon thường nói, khiến tôi không khỏi khẽ cười lớn. Cô ấy cũng đã học được cách nói chuyện giống anh trai mình từ lúc nào không hay.

``Em cứ yên tâm làm việc của em đi. Bọn chị ở nhà sẽ lo liệu mọi thứ, không có gì phải lo lắng đâu. Anh ấy sẽ ổn thôi.''

Sau đó Kson chào tạm biệt rồi cúp máy, tiếng bíp nhỏ của điện thoại vang lên khi cuộc gọi kết thúc. Tôi đặt chiếc điện thoại xuống bàn trà, lúc này Ryuuhou đang ngồi đối diện tôi cũng ngẩng đầu lên, ánh mắt có chút ngưỡng mộ.

``Kson-nee thật sự tự mình xử lý được hết mọi việc ở kho đó chị nhỉ. Em không nghĩ là cô ấy có thể làm được tất cả một mình như vậy.''

``Ừ. Chị ấy đã có nhiều kinh nghiệm làm việc cùng anh Kuon từ lâu rồi, nên cũng nắm rõ hết các quy trình.'' tôi đáp lại, tay vân vê mép gối sofa vô thức.

``Nhưng hôm nay thật sự khác. Bình thường đều có Onii ở bên cạnh để nhắc nhở từng bước, giải đáp mọi thắc mắc. Hôm nay cô ấy phải tự mình đưa ra tất cả các quyết định.'' Ryuuhou nói nhỏ, như thể cũng nhận ra sự thay đổi lớn lao đó.

Tôi nhìn vào màn hình điện thoại đã tắt đen, lòng tràn đầy niềm tự hào không nói nên lời. Tôi đáp lại em ấy, giọng chắc chắn: ``Đúng vậy. Hôm nay chị ấy đã thực sự tự đứng ra một mình, và cô ấy đã làm rất tốt.''

\ct

Khi ánh nắng cuối buổi chiều vừa rủ xuống qua khe cửa sổ tầng năm, Kuon dần thức giấc sau một giấc ngủ dài.

Anh từ từ ngồi dậy, từng cử động chậm rãi như thể mọi sức lực đều bị cơn sốt rút cạn. Mái tóc anh rối tung vì nằm lâu ngày, và khi anh mở lời hỏi thăm, giọng nói vẫn còn khàn đặc, cổ họng dường như vẫn chưa hết rát rỉ vì cơn ho hồi sáng.

Tôi có mặt ở khung cửa phòng gần như ngay lập tức—thực ra từ lúc nghe tiếng chuyển mình đầu tiên vang ra từ trong phòng, tôi đã đứng dậy khỏi chiếc ghế gỗ bên ngoài, bước chân nhẹ nhàng đến bên cạnh giường trước khi anh kịp định thần.

``Anh thấy thế nào rồi?'' tôi hỏi nhỏ, giọng vẫn giữ lại chút thận trọng như sợ làm anh giật mình.

``Đỡ hơn lúc sáng nhiều,'' anh đáp, đầu hơi nghiêng về phía tôi, đôi mắt vẫn còn mơ màng vì cơn mệt mỏi chưa tan hết.

``Có thấy chóng mặt không?'' tôi lại hỏi, mắt dán vào gương mặt anh để quan sát từng biểu cảm nhỏ nhất.

``Một chút, ngay khi vừa ngồi dậy thôi,'' anh thừa nhận, tay khẽ vịn vào mép giường để giữ thăng bằng.

``Vậy thì nằm xuống lại đi,'' tôi lập tức nói, bước đến gần để đỡ lưng anh nếu cần. ``Còn chưa khỏe mạnh gì mà cố gắng ngồi dậy làm gì.''

Kuon ngước nhìn tôi, khóe miệng khẽ nhếch lên một chút như thể muốn phản đối. ``Anh chỉ mới ngồi dậy có mấy giây thôi mà.''

``Và anh vừa nói mình đang bị chóng mặt mà,'' tôi nhấn mạnh, không chịu nhượng bộ chút nào. ``Nghe lời, nằm xuống nghỉ thêm đi.''

Ngay lúc đó, chiếc đồng hồ thông minh trên bàn cũng bật lên tiếng của Game, giọng điện tử đều đặn như mọi khi: ``Kuon-user nên tiếp tục duy trì nghỉ ngơi tuyệt đối, tuân thủ đúng theo hướng dẫn của phòng khám nhé. Còn ít nhất mười hai tiếng nữa cậu mới được ngồi dậy lâu dài.''

Tôi bước đến bàn, rót một cốc nước ấm từ chiếc bình giữ nhiệt rồi đưa cho anh. Kuon nhận lấy, nâng cốc lên nhấp nhẹ vài ngụm để làm ẩm cổ họng khô rát.

Sau khi đặt cốc nước xuống, anh lại ngước nhìn tôi, bắt đầu lần lượt hỏi thăm những công việc anh để lại: ``Kson đã đi đến kho rồi à? Em có nghe tin gì từ cô ấy không?''

``Ừ, chị ấy đã ở đó từ sáng sớm. Mọi việc ở kho đều ổn, cô ấy tự mình xử lý được hết mọi thứ không cần hỏi anh gì cả,'' tôi đáp.

``Còn Ryuuhou, em ấy đã nhận bộ lịch họp anh để lại chưa? Tất cả các cuộc họp hôm nay đều cần dời sang ngày mai.''

``Em ấy đã xử lý xong hết rồi, tất cả những việc gấp đều được chuyển giao cho người phụ trách dự phòng, không có gì bị sót lại đâu.''

``Suzuno vẫn đang ở bếp nấu cháo phải không? Nồi cháo sáng nay còn không?''

``Anh đã ăn một phần nhỏ lúc chiều rồi mà, quên rồi à?'' tôi mỉm cười, nhắc lại chuyện hồi chiều Suzuno mang cháo lên. ``Còn lại đều được cất vào hộp giữ nhiệt trong tủ lạnh, anh đói thì mai hâm lại ăn cũng được.''

``Còn Naomi thì sao? Em bé ấy vẫn ngoan chứ?''

Tôi chỉ tay vào bức tranh nhỏ của Naomi, vẫn được đặt dựa vào chiếc cốc nước trên bàn cạnh giường, nơi mà anh có thể nhìn thấy ngay khi ngước mắt. Kuon hướng mắt theo hướng tôi chỉ, ánh mắt vốn còn mệt mỏi bỗng dịu đi hẳn, như thể mọi lo âu trong người đều tan biến khi nhìn thấy bức vẽ mặt trời rực rỡ của em bé.

``Mọi người trong nhà đã làm quá nhiều chuyện vì anh rồi,'' anh khẽ nói, giọng có chút xót xa. ``Từ sáng đến giờ mọi người đều lo liệu hết mọi thứ, anh chỉ nằm đây mà không giúp được gì.''

``Mọi người chỉ làm phần việc của mình thôi mà,'' tôi đáp lại nhẹ nhàng, ``chẳng có ai làm quá sức đâu, anh đừng lo lắng về chuyện đó.''

Kuon dựa lưng vào chiếc gối dày phía sau, đã từ từ nằm xuống theo lời tôi đề nghị. Anh nhìn tôi đứng bên cạnh giường, rồi nói: ``Vậy em cũng đi nghỉ đi nhé, từ sáng đến giờ chị cũng chẳng nghỉ ngơi được bao lâu. Đứng đây canh anh mãi cũng không phải cách.''

Tôi lắc đầu, không đồng ý với lời đề nghị của anh. ``Chị sẽ nghỉ sau khi kiểm tra lại nhiệt độ cho anh lần này nữa.''

``Em vừa kiểm tra nhiệt độ cho anh mười lăm phút trước mà,'' anh nhắc lại, nhớ rõ khoảng thời gian tôi vừa đo nhiệt độ hồi nãy.

``Chị biết,'' tôi vẫn giữ nguyên lời nói của mình, ``nhưng chị muốn kiểm tra lại thêm lần nữa để chắc chắn.''

Tôi đứng yên bên cạnh giường, tay vô thức cử động như muốn chạm lên trán anh để cảm nhận nhiệt độ, nhưng lại cố gắng kìm lại—tôi đã làm vậy quá nhiều lần trong ngày, sợ rằng bản thân sẽ làm phiền đến giấc ngủ của anh. Ngay lúc tôi đang phân vân, Game đã đọc to con số nhiệt độ vừa đo xong từ chiếc nhiệt kế trên bàn.

Nhiệt độ của anh đã giảm thêm một chút so với lần kiểm tra trước, không còn giữ ở mức cao ngất ngưởng như hồi sáng nữa.

Tôi thở ra một hơi thật khẽ, hơi thở ấy nhẹ đến mức chính tôi cũng gần như không cảm nhận được, nhưng tất cả sự nặng nề trong lòng từ sáng đến giờ bỗng dưng nhẹ bớt đi phần nào, như thể một tảng đá lớn vừa được nhấc đi.

\pov{Kson}

Tôi bước chân vào cổng kho đúng tám giờ rưỡi sáng, hơi lạnh của buổi sớm còn vương trên áo khoác. Ánh mặt trời mới nhú lên chiếu rọi vào những chồng thùng hàng xếp gọn gàng, tiếng động cơ xe tải đầu ngày bắt đầu vang lên từ xa. Bà Kaien đã đợi sẵn trong văn phòng nhỏ góc kho, trên bàn đặt một tờ giấy in chi tiết danh sách tất cả công việc cần hoàn thành trong ngày hôm nay, được đánh dấu màu xanh lá cây cho những việc ưu tiên, đỏ cho các đơn hàng cần xử lý gấp.

Bao lâu nay, mỗi lần đến kho, anh Kuon luôn là người có mặt sớm hơn tôi. Anh thường đứng trước tấm bảng lịch giao hàng treo trên tường, mắt lướt qua từng địa điểm, từng khung giờ, rồi cầm bút sắp xếp thứ tự các cuộc gọi cần thực hiện, phân chia công việc rõ ràng cho từng người trong đội. 

Ngày hôm nay, chiếc ghế gỗ rộng rãi ở đầu bàn mà anh vẫn ngồi mỗi ngày để trống, mặt ghế còn lưu lại chút vết bụi mỏng vì không có ai ngồi lên. Tôi đứng nhìn chiếc ghế đó vài giây, nghe tiếng gió thổi qua cửa sổ, rồi mới từ từ kéo ghế ra, ngồi vào vị trí quen thuộc mà trước đây chỉ có anh mới đảm nhận.

``Cháu muốn bắt đầu từ đâu trước? Danh sách giao hàng hay các cuộc gọi cần xác nhận với nhà cung cấp?'' bà hỏi, giọng nhẹ nhàng như thường lệ, tay chỉ vào tờ lịch chi tiết trên bàn. Bà đưa cho tôi cuốn sổ lịch giao hàng bìa cứng, các trang giấy đã hơi cũ vì được lật đi lật lại nhiều lần.

``Cháu xem đơn nào cần xử lý trước. Những việc thông thường cứ làm theo đúng thứ tự trong danh sách này. Nếu gặp vấn đề gì vượt quá quyền hạn mà cháu không thể tự quyết định, thì gọi cho cô.''

``Rõ ạ,'' tôi đáp lời, cầm lấy cây bút bi đen trên bàn. Mở cuốn sổ, tôi lướt mắt qua từng đơn hàng, rồi dùng bút đánh dấu gạch chân ba đơn có giờ giao hàng sát nhau chỉ cách nhau hai mươi phút, cần phải sắp xếp lộ trình cẩn thận để không bị trễ.

Vừa lúc tôi đang tính toán lộ trình, một nhân viên phụ trách bảo trì kho bước vào văn phòng, báo rằng xe tải chạy tuyến phía đông đột ngột bị hỏng lốp, cần phải kéo vào gara sửa chữa, ít nhất cũng mất một tiếng rưỡi mới có thể hoạt động lại. 

Tôi đứng dậy bước đến tấm bản đồ lớn treo trên tường, nơi đánh dấu tất cả các điểm giao hàng trong thành phố, ngón tay chỉ lần lượt qua hai địa điểm. Sau vài phút suy nghĩ, tôi quyết định đổi thứ tự giao hàng của hai đơn còn lại, sắp xếp cho xe tải vẫn hoạt động bình thường đi vòng qua các điểm đó, đảm bảo chúng vẫn đến đúng giờ hẹn với khách hàng.

Bà Kaien quan sát kỹ lộ trình tôi vừa vạch ra trên bản đồ, rồi khẽ gật đầu như để công nhận: ``Cháu cũng suy nghĩ thấu đáo lắm đấy chứ.''

Tôi quay lại bàn làm việc, lấy điện thoại gọi cho tài xế của xe còn lại, nói rõ từng địa chỉ mới và khung giờ dự kiến sẽ đến, giải thích rõ nguyên nhân thay đổi lộ trình. Người tài xế lắng nghe và đồng ý thực hiện, không có bất kỳ phàn nàn nào. Cúp máy, tôi đứng yên một lúc, nhận ra rằng trong suốt quá trình xử lý sự cố, không ai phải gọi điện cho Kuon để hỏi ý kiến. Không cần anh phải can thiệp, mọi việc vẫn được giải quyết ổn thỏa.

Tôi nhìn chiếc điện thoại trong tay, màn hình vẫn sáng lên cuộc gọi vừa kết thúc. Một cảm giác lạ lẫm tràn vào lòng: Tôi thực sự đã tự mình làm được. Chuyện đó không phải là phép màu gì, chỉ là một quyết định nhỏ bé, được suy nghĩ kỹ lưỡng và kiểm tra với người có kinh nghiệm trước khi thực hiện. Nhưng sáng nay, cái quyết định nhỏ bé ấy lại mang ý nghĩa lớn hơn tôi từng tưởng tượng, như thể tôi vừa bước qua một ngưỡng cửa mới trong công việc.

Tôi quay lại vị trí ngồi trên ghế của anh Tổng, hỏi mẹ của anh ấy: ``Cô còn có đơn nào khác cần em xem xét xử lý không ạ?''

Mẹ Kuon nở một nụ cười ấm áp, đôi mắt ánh lên sự tự hào: ``Có đấy. Để cô đưa cuốn sổ ghi nhận tồn kho cho cháu.''

Tôi nhận lấy cuốn sổ bìa nâu từ tay bà, lật mở trang đầu tiên. Ngày làm việc chính thức của tôi tại kho, khi phải tự mình đứng ra chịu trách nhiệm, bắt đầu từ đây.

\ct

Đến gần giữa trưa, danh sách công việc dài dằng dặc trên bàn đã ngắn đi một nửa, những mục đã hoàn thành được tôi gạch bỏ bằng bút xanh. Tôi kiểm tra lại số lượng vật tư còn lại trong kho, ký nhận hai kiện hàng mới vừa được giao đến từ nhà cung cấp, và gọi điện xác nhận với khách hàng về một đơn hàng bị trễ do thời tiết, giải thích rõ tình hình và hẹn lại thời gian giao hàng mới.

Có một lần, khi gặp phải khoản hoàn hàng từ một khách hàng hủy đơn, tôi vô tình mở cuộc trò chuyện với anh KUon trên điện thoại, ngón tay đang định gõ câu hỏi về cách ghi nhận khoản tiền đó vào sổ sách. Nhưng rồi tôi dừng lại, nhìn vào ô nhập tin nhắn trống trắng. Tôi nhớ ra anh đang nằm nghỉ ở nhà vì bị ốm, không thể làm phiền anh với những công việc ở kho. 

Thay vào đó, tôi bước ra tìm bà Kaien, hỏi bà về quy trình xử lý. Bà chỉ cho tôi biểu mẫu cần điền, giải thích từng mục rõ ràng. Tôi tự mình điền đầy đủ thông tin, lưu lại phiếu vào hệ thống máy tính. Thực ra chẳng có gì phức tạp cả, chỉ là tôi đã quá quen với việc có anh Kuon ở bên cạnh, sẵn sàng trả lời bất cứ câu hỏi nào tôi đưa ra ngay khi tôi còn phân vân.

Ngay lúc đó, chiếc đồng hồ thông minh của Game trên bàn bật sáng, gửi một tin nhắn vào nhóm trò chuyện gia đình: ``Kuon-user đang trong thời gian nghỉ ngơi. Xin đừng gửi bất kỳ câu hỏi liên quan đến công việc tới thiết bị của cậu ấy.''

Tôi nhanh chóng nhấn trả lời: ``Biết rồi. Hôm nay tôi sẽ tự mình xử lý được hết mọi việc ở kho.''

Game phản hồi lại chỉ bằng một biểu tượng dấu kiểm xanh nhỏ, xác nhận đã nhận được tin nhắn của tôi. Tôi tắt màn hình điện thoại, cất nó vào túi quần.

Một nhân viên kho bước vào văn phòng, mang theo một tập hóa đơn mới vừa được gửi đến. Tôi ngồi xuống kiểm tra từng số đơn hàng, ghép từng hóa đơn vào đúng bìa hồ sơ của khách hàng đó, đảm bảo không có nhầm lẫn hay sót sót gì. Sau khi tôi hoàn thành, mẹ Kuon xem lại toàn bộ công việc của tôi, rồi ngẩng đầu nói: ``Cháu làm nhanh hơn cô nghĩ đấy. Mọi thứ đều rất chính xác.''

``Dạ, cũng vì Kuon-paisen đã chỉ dẫn em cách làm nhiều lần rồi ạ.'' tôi đáp lại một cách tự nhiên, vẫn đang sắp xếp các bìa hồ sơ vào kệ.

``Biết làm theo đúng hướng dẫn của người khác là một chuyện. Tự mình đứng ra chịu trách nhiệm, tự mình đưa ra quyết định và hoàn thành công việc lại là một chuyện hoàn toàn khác,'' bà ấy nói nhẹ nhàng.

Tôi không biết phải đáp lại câu nói đó như thế nào, chỉ có thể mỉm cười, tiếp tục xếp các hóa đơn vào đúng vị trí. Bên ngoài cửa kho, những đám mây xám kéo dần qua bầu trời, có vẻ như sắp có mưa, nhưng cơn mưa vẫn chưa rơi xuống. Tôi quay lại với danh sách công việc còn lại, tiếp tục hoàn thành từng mục một, cảm thấy tự tin hơn vào khả năng của mình.

\ct

Khi đồng hồ điểm mười hai giờ trưa, tôi nhấc hộp cơm nguội lạnh đặt từ sáng lên, bước ra bậc thềm bê tông phía sau kho để nghỉ ngơi. Gió thổi qua mấy bụi cỏ dại mọc ven tường, mang theo chút mùi đất ẩm từ lúc sớm mai, cuối cùng mới có vài phút để thở phào sau buổi sáng bận rộn chạy qua chạy lại giữa các đơn hàng. 

Tôi ngồi trên bậc đá mát, mở hộp cơm ra rồi lướt danh bạ điện thoại, nhấn gọi cho chị Kaigou, bởi cả sáng nay tôi vừa làm việc vừa lo lắng không yên về tình hình của anh Kuon ở nhà.

Chỉ có hai hồi chuông, chị đã bắt máy ngay, tiếng đáp vọng nhẹ từ đầu dây. ``Anh ấy đang ngủ. Em có việc gấp à?''

``Không có gì gấp cả đâu,'' tôi vội trả lời, ngón tay vô thức xoắn lấy mép hộp cơm. ``Em chỉ muốn hỏi là anh ấy có uống đủ nước không, có ai nhớ mang nước lên cho anh ấy không?''

``Không phải lo. Suzu đã lập bảng giờ uống nước, cứ mỗi tiếng lại mang lên một cốc ấm cho anh ấy, không bao giờ quên đâu,'' Kaigou nói tiếp, giọng có chút mệt mỏi như vừa chạy đi chạy lại cả buổi. ``Còn gì nữa không?''

``Nhiệt độ của anh ấy thì sao? Có hạ được chút nào không ạ?'' tôi hỏi nhanh, tim tự nhiên đánh nhanh hơn khi chờ câu trả lời.

``Vẫn còn cao, nhưng đã giảm được nửa độ so với lúc sáng rồi. Có tiến triển là được,'' chị đáp, và tôi có thể nghe ra giọng chị căng thẳng hơn nhiều so với ngày thường, như thể vẫn đang nín thở chờ đợi từng con số nhiệt độ mỗi lần đo.

Tôi do dự một chút rồi hỏi nhỏ: ``Chị ở nhà đang ổn chứ ạ? Có cần em về giúp gì không?''

``Chị ổn, mọi thứ đều trong tầm kiểm soát,'' Kaigou trả lời liền, nhưng giọng vẫn không thỏa được cho lắm.

``Chị chắc chắn không có gì sao ạ? Nếu có việc gì--''

``Kson,'' chị gọi tên tôi, ngắt lời nhẹ nhàng như muốn bảo tôi đừng lo lắng thêm.

``\dots Rồi ạ, em không hỏi nữa. Em sẽ cố gắng xử lý hết mọi việc ở kho, chị đừng lo về công việc,'' tôi vội đáp, hiểu được ý chị muốn nói.

Đầu dây im lặng một lúc, rồi Kaigou nói nhỏ hơn, như chỉ sợ mình nói to sẽ làm phiền ai đó trong nhà: ``Thật sự cảm ơn em đã nhận ra đến kho thay anh ấy hôm nay. Nếu không có em, mẹ chắc đã phải xoay sở vất vả lắm.''

``Em tự nguyện mà chị, đâu có gì. Anh ấy cũng luôn giúp đỡ em từ lúc mới đến làm mà,'' tôi cười, cảm thấy lòng ấm áp lên khi nghe chị nói vậy.

``Chị biết. Nhưng vẫn phải cảm ơn. Em ăn trưa đi nhé, có gì gọi lại sau.''

Cuộc gọi kết thúc, tiếng bíp nhỏ vang lên. Tôi đặt điện thoại xuống bên hộp cơm, dạo chút cơm vào miệng. Lần đầu tiên kể từ khi bắt đầu ca làm sáng nay, tôi không cảm thấy thôi thúc phải nhắn tin hay gọi cho Kuon để xin xác nhận bất cứ điều gì. Tất cả những gì cần hỏi, tôi đã hỏi đúng người, và mọi thứ vẫn đang diễn ra theo đúng lộ trình, không có gì bị kẹt hay trì hoãn. Tôi ăn hết phần cơm còn lại trong hộp, lau miệng rồi đứng dậy, bước quay vào trong kho để tiếp tục công việc còn lại.

\ct

Vừa mới bắt đầu giờ chiều, một thông báo đột ngột hiện lên trên hệ thống: một khách hàng yêu cầu đổi địa chỉ nhận hàng cho đơn hàng số bảy, từ địa chỉ cũ ở quận ba chuyển sang một văn phòng mới ở quận năm. Tôi mở đơn hàng ra kiểm tra lại toàn bộ thông tin, ghi lại địa chỉ mới cẩn thận, rồi lập tức gọi lại cho khách hàng để xác nhận bằng lời, tránh nhầm lẫn hay sai sót gì có thể xảy ra. 

Sau khi khách hàng xác nhận đúng địa chỉ mới, tôi cập nhật lại thông tin lên hệ thống, gọi cho tài xế phụ trách tuyến đó để thông báo thay đổi lộ trình, đảm bảo anh ấy vẫn đến đúng giờ và giao hàng đúng người. Tất cả các bước tôi đều làm theo đúng quy trình anh Kuon đã dạy từ trước, không có bước nào bị sót hay làm sai.

Nhưng đến khi cần gửi báo cáo cuối ngày, tôi dừng tay lại trước nút ``gửi'' trên màn hình máy tính. Ngón tay treo lơ trên chuột, một cảm giác quen thuộc dâng lên: tôi muốn gửi bản báo cáo này cho anh ấy xem qua trước, muốn nghe anh ấy gật đầu nói ``được rồi'' mới dám bấm gửi. Nhưng rồi tôi chợt nhớ ra, anh ấy đang nằm nghỉ ở nhà vì cơn sốt, không thể bị làm phiền bởi những công việc ở kho này.

Tôi thở nhẹ một hơi, rồi kéo chuột đọc lại từng dòng trong báo cáo, kiểm tra lại từng con số, từng địa chỉ, từng mã đơn hàng một lần nữa. Sau đó tôi gọi bà Kaien lại xem qua, bà đọc lướt qua các mục chính rồi cũng gật đầu xác nhận không có gì sai sót. Lúc đó tôi mới nhấn nút gửi.

Thông báo ``đã gửi thành công'' hiện lên trên màn hình, tôi thở phào nhẹ nhõm rồi nói với bà: ``Xong rồi ạ.''

Bà Kaien dựa vào mép bàn, gật đầu nhìn tôi với nụ cười ấm áp: ``Hôm nay cháu đã thay Kuon xử lý được gần như tất cả các việc thường nhật ở kho rồi đấy. Mọi thứ đều ngăn nắp, đúng quy trình cả rồi.''

``Trước đây cháu cứ nghĩ mình sẽ phải gọi hỏi anh ấy suốt ngày, không thể tự mình xử lý hết được,'' tôi thừa nhận, sắp xếp các giấy tờ đã ký vào tập hồ sơ. ``Nhưng hôm nay lại làm được hết, cháu cũng không ngờ tới ạ.''

``Hỏi khi cần hỗ trợ không có gì xấu đâu, mà hôm nay cháu đã tự mình tìm được người có thể xác nhận, tự mình đứng ra đưa ra quyết định\dots\ đó mới là điều quan trọng,'' bà nói nhẹ nhàng, tay vỗ nhẹ vào vai tôi.

Tôi vô thức quay đầu nhìn về chiếc ghế gỗ rộng ở đầu bàn, nơi mà Kuon vẫn ngồi mỗi ngày. Chiếc ghế vẫn trống, nhưng trong lòng tôi đã không còn cảm giác bỡ ngỡ, lạ lẫm khi ngồi vào vị trí đó nữa. ``Em cũng không biết mình lại cần nghe câu đó đến thế. Nghe bà nói vậy, em thấy tự tin hơn hẳn.''

``Giờ thì cháu biết rồi, phải không?'' bà cười.

Tôi cũng cười theo, đứng dậy cất tập hồ sơ vào kệ trên tường. Ca làm việc vẫn chưa kết thúc, vẫn còn vài việc nhỏ cần xử lý trước khi đóng cửa kho, nhưng tôi biết phần khó nhất, phần phải đối mặt với mọi tình huống một mình mà không có Kuon ở bên cạnh, đã qua rồi.

\ct

Khi tất cả các công việc cuối cùng trong ngày đã được xếp gọn vào kệ hồ sơ, tất cả các xe tải đã về đến bến và cửa kho được khóa lại an toàn, tôi rút điện thoại ra gọi về nhà một lần cuối trước khi bước lên xe cùng bà Kaien để về nhà. Tôi đã chuẩn bị sẵn danh sách tất cả những việc đã hoàn thành để báo cáo, chắc chắn rằng không có gì bị sót sót hay trì hoãn.

Lần này, không phải chị Kaigou hay ai khác nhấc máy mà chính Kuon. Tiếng nhấc máy vang lên, và ngay lập tức giọng anh ấy vang qua đầu dây: ``Kson à? Có chuyện gì phát sinh ở kho không em?''

Tôi có thể nghe rõ từng chữ mà anh ấy cố gắng nói một cách bình thường, như thể muốn che giấu tình trạng sức khỏe của mình. Nhưng tiếng khàn khốc, rát rỉ vì cơn sốt và ho vẫn lộ ra rõ ràng trong từng âm tiết, làm trái tim tôi vô thức thắt lại.

``Không có chuyện gì đâu ạ,'' tôi nhanh chóng nói, để anh ấy khỏi lo lắng. ``Em chỉ gọi để báo cáo thôi, mọi thứ ở kho đều ổn cả.''

``Em đã xử lý hết hết các việc mà anh để lại chưa?'' giọng anh ấy vẫn cố gắng giữ ổn định, nhưng có thể nghe ra sự mệt mỏi ngấm vào từng lời nói.

``Đơn hàng yêu cầu đổi địa chỉ đã được em xác nhận lại với khách hàng rồi. Lịch giao hàng cho ngày mai đã cập nhật hoàn chỉnh lên hệ thống, tất cả các tài xế đều được thông báo lộ trình mới. Còn báo cáo cuối ngày của hôm nay cũng đã gửi cho nhà cung cấp đúng giờ rồi.''

``Còn khoản hoàn hàng từ khách hàng hủy đơn hôm nay thì sao? Em xử lý thế nào rồi?'' anh ấy tiếp tục hỏi, như thể muốn kiểm tra tất cả các khâu cuối cùng.

``Em đã dùng đúng biểu mẫu ghi nhận trong sổ mà anh ấy vẫn dùng, bà Kaien đã kiểm tra lại phiếu đó và ký xác nhận rồi, không có sai sót gì đâu ạ.''

Anh Kuon im lặng một nhịp ngắn, như thể đang nghỉ ngơi sau khi hỏi hết các câu hỏi, hay đang trong quá trình xử lý thông tin tôi vừa đưa ra. Rồi anh ấy nói nhẹ nhàng: ``Tốt lắm. Vậy thì nhớ--''

Chưa kịp nói hết câu, anh ấy đột ngột ho sặc sụa. Tiếng ho vang qua điện thoại, kéo dài liên tục, đến mức tôi đứng thẳng dậy khỏi ghế, tay vô thức siết chặt chiếc điện thoại, lo lắng không thôi.

``Kuon-paisen? Anh ấy ổn chứ?'' tôi gọi vội, giọng có chút hoảng sợ.

Ngay lập tức, có tiếng bước chân vội vã từ đầu dây bên kia, rồi chị Kaigou giành lấy chiếc điện thoại từ tay anh Kuon, giọng chị có chút bực bội nhưng đầy lo lắng: ``Aniki, em đã bảo anh bao nhiêu lần là không được nghe máy khi còn ốm. Sao anh lại không nghe lời?''

``Anh chỉ muốn hỏi Kson xem công việc ở kho thế nào thôi mà--'' anh Kuon cố gắng giải thích, nhưng câu nói bị ngắt quãng bởi một cơn ho khác, còn mạnh hơn lúc trước.

Tôi nhanh chóng lấy lại bình tĩnh, cố gắng giữ giọng của mình thật nghiêm nghị, như cách mà anh ấy vẫn thường làm với mọi người khi cần thiết: ``Aniki, cúp máy ngay đi rồi uống một ngụm nước ấm đi ạ. Tất cả báo cáo công việc đã xong hết rồi, không còn gì nữa mà cần anh ấy phải chỉ đạo hay can thiệp đâu.''

``Nhưng mà--'' anh ấy vẫn còn muốn nói gì đó, nhưng tôi đã ngắt lời.

``Không có nhưng gì hết ạ. Em đã tự mình xử lý được hết mọi việc ở kho rồi, không có gì trở ngại cả. Anh ấy cứ yên tâm nghỉ ngơi đi, đừng nghĩ đến công việc nữa.''

Chị Kaigou cũng lập tức nói vọng vào, đồng tình với lời tôi: ``Nghe lời Kson-chan đi anh, cứ nghỉ ngơi đi thôi.''

Một tiếng thở dài yếu ớt vang lên từ đầu dây, tiếng thở dài mệt mỏi của anh Kuon. Cuối cùng anh ấy cũng chịu thua: ``Được rồi. Anh cúp máy đây. Các em cũng nhớ nghỉ ngơi sớm.''

Cuộc gọi kết thúc, màn hình điện thoại từ từ tối đi, để lại phản chiếu khuôn mặt của tôi trên nền đen. Tôi đứng yên một lúc, vẫn còn cảm thấy lo lắng về tình hình của anh Kuon ở nhà.

Bà Kaien đã đứng gần cửa văn phòng kho từ lúc nào, tay cầm chìa khóa xe của mình, nhìn tôi với một nụ cười ấm áp. Bà nói nhỏ: ``Vất vả cho cháu ngày hôm nay rồi.''

Tôi cất chiếc điện thoại vào túi quần, rồi đi theo bà ra khỏi văn phòng, khóa cửa lại sau lưng. Tôi đáp lại nhẹ nhàng: ``Hôm nay con chỉ làm đúng phần việc của mình thôi mà, đâu có làm gì quá đáng đâu ạ.''

Bà dừng bước một lát, rồi đặt tay nhẹ lên vai tôi, giọng nói trầm ấm: ``Đôi khi, chỉ cần làm tốt phần việc của mình, đứng ra chịu trách nhiệm khi cần thiết, là đã đủ rồi.''

\pov{Mikeneko}

Ánh hoàng hôn cuối cùng cũng tàn đi sau những ngọn nhà xa xôi, và căn nhà Hikawa dần chìm vào im lặng khi đêm buông xuống từng tầng một. Ánh đèn vàng ấm của phòng khách ở tầng năm đã tắt từ lúc nào, thay vào đó chỉ còn ánh sáng nhợt nhạt từ mặt đường lọt qua khe cửa sổ. Tiếng ồn ào của bộ phim truyền hình mà mọi người cùng xem buổi tối đã ngừng hẳn, không còn tiếng cười nói vang vọng giữa các căn phòng nữa. 

Naomi và Koneko đã được đưa lên tầng tám từ hơn một tiếng trước, giờ chắc hẳn hai đứa trẻ đã chìm vào giấc ngủ say, tròn mắt trên những chiếc gối mềm. Kson vừa về đến nhà sau ca làm việc kéo dài cả ngày ở kho, chiếc áo khoác của em ấy vẫn còn vương chút hơi lạnh của đêm sớm, còn Delutaya cũng vừa đóng cửa phòng thu, khóa lại những thiết bị âm thanh sau một ngày dài làm việc.

Tất cả mọi người đều đã ổn định, chỉ có tôi vẫn còn ngồi trên giường, mắt mở lớn nhìn vào bóng tối của căn phòng mình. Không phải vì có việc gì gấp cần giải quyết, cũng không phải vì tôi còn phải làm việc gì chưa xong. Chỉ là từ sáng đến tối, tôi đã nghe thấy tiếng bước chân của mọi người qua lại cầu thang, tiếng gọi hỏi nhau về tình trạng của anh Kuon, rồi dần dần những tiếng động ấy nhỏ lại, cho đến khi cả căn nhà chỉ còn tiếng tích tắc của đồng hồ treo tường ở hành lang.

Từ sáng sớm, tôi đã biết anh Tổng bị ốm. Chị Kaigou đã nhắn tin vào nhóm trò chuyện chung của gia đình ngay khi phát hiện anh ấy sốt cao, tin nhắn đầu tiên của chị ấy gửi đi lúc bảy giờ sáng. Tôi đã mở tin nhắn đó ngay khi nhận được, nhưng ngón tay treo lơ trên bàn phím, không biết phải gõ gì để trả lời. Tôi không quen với việc thể hiện sự lo lắng ra ngoài, không biết nên nói gì để tỏ ra tôi cũng quan tâm, như những người khác trong nhà.

``Anh ấy sốt cao, đang nằm nghỉ ở tầng năm. Mọi người đừng lên đó làm ồn ảnh hưởng đến anh ấy nghỉ ngơi,'' là tin nhắn đầu tiên của Kaigou trong nhóm. Sau đó, Suzuno gửi một bức ảnh bát cháo nóng hổi vừa nấu xong, ghi chú là chuẩn bị mang lên cho anh ấy.

Naomi gửi bức vẽ màu rực rỡ của em bé, hình ảnh một mặt trời vàng ươm với những tia sáng tỏa ra xung quanh, kèm theo những chữ cái xiên xẹo mà chỉ người thân mới đọc được: ``Chúc Onii-chan mau khỏe.'' Ryuuhou cập nhật tin tức về công việc, rằng tất cả lịch trình và cuộc họp của Kuon đã được chuyển giao hết, không có việc gì bị trì hoãn.

Tôi đã nhìn lướt qua từng tin nhắn đó, lật đi lật lại trong cuộc trò chuyện của nhóm hơn mười lần. Màn hình điện thoại sáng lên trong bóng tối của căn phòng, chiếu lên mặt tôi một vệt sáng mỏng. Cuối cùng tôi cũng không nhắn lại gì cả, chỉ đặt chiếc điện thoại xuống chiếc bàn cạnh giường, để màn hình tự tối đi. 

Tôi tự nhủ trong lòng rằng cả nhà đã ở bên cạnh anh ấy rồi, chị Kaigou thì luôn đứng trước cửa phòng anh ấy để chăm sóc, Suzuno thì chuẩn bị đồ ăn nóng hổi, mọi người đều đã lo liệu hết mọi thứ. Tôi không cần phải xuống tầng năm, không cần phải thêm vào một sự lo lắng nữa, mọi thứ đã đủ ổn thỏa rồi.

Nhưng thế nào đi nữa, trong đầu tôi vẫn cứ canh cánh một điều gì đó, như thể có một ngọn đèn nhỏ chưa tắt hẳn, khiến tôi không thể nào chìm vào giấc ngủ được. Căn phòng trong tâm trí tôi vẫn còn sáng lên một chút, không chịu tối hẳn như những gì tôi tự nhủ.

\ct

Khi đồng hồ trong nhà chỉ gần đến nửa đêm, toàn bộ căn nhà đã chìm vào im lặng hoàn toàn, tôi mới nhẹ nhàng nhấc mình khỏi giường. Hành lang tầng tám lúc này im lìm đến mức có thể nghe thấy tiếng gió thổi qua khe cửa sổ ở cuối hành lang, không có một tiếng động nào khác. Tôi khoác vội chiếc áo len dày của mình lên người, đi chân không trên sàn nhà gỗ mát để bước chân được nhẹ nhất có thể, không muốn làm phiền đến ai đang ngủ.

Trên hai tay tôi cầm một chiếc cốc sứ nhỏ, đã được rót đầy nước ấm từ bình giữ nhiệt trong phòng mình. Trong túi áo len, tôi giấu một viên thuốc hạ sốt còn nguyên trong vỉ, đúng loại mà Kuon vẫn hay dùng khi bị ốm, loại thuốc mà gia đình luôn để sẵn trong ngăn thuốc chung ở tầng một. Trước khi lấy ra, tôi đã đọc kỹ nhãn thuốc, kiểm tra lại hướng dẫn sử dụng, liều lượng và cả ngày hết hạn, đảm bảo không có gì sai sót.

Tôi không hề có ý định tự ý đánh thức anh ấy dậy, không muốn làm phiền đến giấc ngủ của anh ấy khi anh ấy đang cần nghỉ ngơi nhất. Tôi chỉ muốn đặt cốc nước đó ở trong phòng, nếu như anh ấy tỉnh dậy vào ban đêm vì cổ họng khô rát, hay cơn sốt làm anh ấy mệt mỏi, thì sẽ có nước ấm để uống, và viên thuốc cùng tờ hướng dẫn sẽ nằm ngay bên cạnh, tiện lợi cho anh ấy sử dụng.

Tôi đi nhẹ nhàng xuống cầu thang, đến trước cửa phòng của anh Kuon ở tầng năm, rồi dùng tay đẩy cánh cửa thật chậm, chỉ để hé một khe nhỏ. Ánh sáng đèn yếu ớt từ hành lang lọt vào, chỉ vẽ một dải sáng mỏng trên sàn gỗ của căn phòng, không đủ để làm phiền đến người đang ngủ. Tôi nhìn vào trong, anh ấy đang nằm ngửa trên giường, chăn đắp đến ngực. Trán anh ấy vẫn còn lấm tấm những hạt mồ hôi nhỏ, do cơn sốt chưa hạ hẳn, chiếc gối bên dưới cũng hơi xê dịch vì anh ấy trở mình vài lần trong đêm.

Trên chiếc bàn nhỏ cạnh giường, chiếc cốc mà mọi người mang lên cho anh ấy ban ngày đã cạn sạch, đặt cạnh chiếc nhiệt kế điện tử vẫn đang bật màn hình nhỏ. Còn có một bát cháo đã được đậy kín nắp, chắc là anh ấy không ăn hết nên để lại, Suzuno có lẽ đã quên mang xuống vào buổi tối. 

Tôi bước nhẹ nhàng đến gần bàn, đặt chiếc cốc nước ấm mà tôi mang theo vào vị trí dễ thấy nhất. Viên thuốc còn nguyên vỉ được tôi đặt cạnh một tờ giấy nhỏ mà tôi đã ghi hướng dẫn sử dụng trước, ghi rõ liều lượng và thời gian cách nhau giữa các lần uống.

Tôi không dám chạm vào người anh ấy, thậm chí không dám phát ra một tiếng động nhỏ nào. Cũng không gọi tên anh ấy để kiểm tra tình trạng, chỉ sợ làm anh ấy giật mình tỉnh giấc. Khi nhìn thấy mép chăn bị tuột xuống, gần như chạm vào sàn nhà, tôi chỉ nhẹ nhàng kéo nó lên một chút, để nó che phủ được chân anh ấy, giữ cho anh ấy không bị lạnh. Sau khi làm xong mọi thứ, tôi lùi dần về phía cửa, giữ nguyên sự im lặng.

\ct

Ngay lúc tôi vừa lùi được vài bước, Kuon đột nhiên khẽ trở mình trên giường, làm cho chiếc gối phát ra một tiếng xào xạc nhỏ. Tôi lập tức đứng bất động tại chỗ, tim đập nhanh lên một chút vì sợ anh ấy tỉnh dậy. Nhưng anh ấy không mở mắt, hơi thở vẫn đều đặn như trước, chỉ có một tiếng ho nhỏ bật ra từ cổ họng, rồi nhanh chóng tắt đi, như thể cơn ho chỉ thoáng qua trong giấc ngủ.

Tôi nhìn lại chiếc cốc nước ấm một lần nữa, chắc chắn rằng nó nằm ở vị trí đủ gần, anh ấy chỉ cần đưa tay ra là có thể với tới khi tỉnh dậy vào lúc nào đó trong đêm. Tờ hướng dẫn sử dụng thuốc cũng nằm ngay bên cạnh, rõ ràng và dễ đọc. Tôi đã làm tất cả những gì mình có thể làm, không có gì còn thiếu sót nữa. 

Tôi quay người, bước nhẹ nhàng về phía cửa, trong đầu vẫn giữ nguyên sự cẩn thận để không làm ra tiếng động. Ngay trước khi khép cánh cửa lại, tôi ngoái đầu nhìn vào trong phòng một lần cuối. Kuon vẫn đang ngủ yên, hơi thở đều đặn, dường như cơn sốt đã bắt đầu dịu đi.

Tôi khép cửa lại gần như không phát ra bất kỳ âm thanh nào, để căn phòng trở lại im lặng như ban đầu. Hành lang bên ngoài cũng nhanh chóng trở lại trạng thái tối mịn như cũ, chỉ còn ánh sáng yếu ớt từ mặt đường. Tôi đi nhẹ nhàng lên cầu thang, trở về căn phòng của mình ở tầng tám, bước chân vẫn giữ được sự nhẹ nhàng nhất. Tôi nghĩ rằng không ai trong nhà đã nhìn thấy tôi ra khỏi phòng, không ai biết tôi đã xuống tầng năm để mang nước và thuốc cho Kuon. Tôi hy vọng vậy, để sự quan tâm thầm lặng này chỉ là một điều nhỏ bé giữa tôi và anh ấy, không cần phải nói ra với ai khác.

\ct

\pov{Hikawa Kuon}

Sáng hôm sau, cơn khô rát ở cổ họng làm tôi thức dậy giữa không gian im lặng của căn phòng. Ánh nắng sớm luồn lách qua khe hẹp của tấm rèm cửa, vẽ nên một vệt sáng mỏng manh chạy dài trên sàn gỗ cũ, như một dải lụa vàng nhẹ nhàng ôm lấy không gian. Đầu tôi vẫn còn cảm thấy nặng nề, mệt mỏi như có thứ gì đó đè lên, nhưng cơn sốt nóng bức đã dịu đi rất nhiều so với đêm trước—không còn cảm giác da mình như bị đốt cháy mỗi khi cử động.

Vẫn nằm trên giường, tôi đưa bàn tay mệt mỏi tìm về chiếc bàn nhỏ cạnh giường, nơi mọi người thường đặt cốc nước cho tôi. Ngón tay chạm vào thành cốc sứ ấm, và tôi ngạc nhiên khi nhận ra nó đã được rót đầy nước ấm, còn nguyên độ ấm vừa đủ để uống. Ngay bên cạnh chiếc cốc, một viên thuốc hạ sốt còn nguyên trong vỉ nhựa nằm gọn, kèm theo một tờ giấy nhỏ ghi rõ liều lượng và cách sử dụng, chữ viết rõ ràng nhưng tôi không thể nhận ra là của ai.

Tôi nhìn chằm chằm vào hai thứ đó, cố gắng lục lọi trong trí nhớ mơ màng của cơn sốt xem có ai đã bước vào phòng đêm qua. Nhưng tất cả đều là mảng mơ hồ, những tiếng động nhỏ trong giấc ngủ không thể thành hình. Tôi chẳng nhớ gì cả, chẳng có một dấu vết nào để lại ngoài chiếc cốc nước đầy và viên thuốc kia.

Đột nhiên, màn hình của chiếc đồng hồ thông minh trên bàn bật sáng, Game xuất hiện với giọng điện tử đều đặn: ``Tôi không ghi nhận có ai bước vào phòng của cậu trong khoảng thời gian từ nửa đêm đến sáng nay.''

Tôi nhíu mắt, hỏi nhỏ: ``Vậy camera hành lang thì sao? Nó không ghi được gì à?''

``Góc quay của camera chỉ bao gồm phần lớn hành lang, không thể chụp được cửa phòng của cậu. Không có dữ liệu nào về người đi qua trước cửa phòng cậu trong đêm qua.''

Tôi lại nhìn chiếc cốc nước thêm một lúc lâu, cảm nhận hơi ấm từ thành cốc lan vào lòng bàn tay. Ai đó trong nhà đã biết tôi sẽ cần nước ấm khi tỉnh dậy, ai đó đã nhớ đến loại thuốc tôi vẫn dùng khi bị ốm, và đã lặng lẽ mang vào mà không muốn ai biết. Tôi không vội gọi mọi người vào phòng để hỏi ngay, không muốn làm phiền không gian bình yên của buổi sáng. Thay vào đó, tôi nâng chiếc cốc lên, uống vài ngụm nước ấm trôi qua cổ họng khô rát, cảm thấy cơ thể nhẹ đi chút ít. Chỉ sau khi đã uống đủ nước, tôi mới nhấc điện thoại, gọi cho Kaigou.

\ct

Kaigou đến phòng của tôi nhanh đến mức tôi nghi rằng em ấy đã đứng ngay bên ngoài cửa, chờ đợi tôi tỉnh dậy từ lâu rồi. Giọng em vừa lo lắng vừa nhẹ nhàng khi bước vào: ``Anh tỉnh rồi à? Cảm thấy thế nào rồi?''

Tôi gật đầu, đưa mắt về chiếc bàn cạnh giường: ``Ừ, đỡ hơn nhiều. Nhưng em có biết ai đã mang nước và thuốc vào đây không? Ai đã vào phòng anh đêm qua?''

Kaigou bước đến cạnh bàn, nhìn chằm chằm vào chiếc cốc nước và viên thuốc, mặt hiện rõ sự ngạc nhiên: ``Không phải em đâu ạ. Em không vào phòng anh sau khi tắt đèn đêm qua.''

Ngay lúc đó, Suzuno vừa mang một chiếc khăn ấm đi qua cửa, nghe thấy câu hỏi của tôi liền dừng bước, lắc đầu: ``Cũng không phải em ạ. Tối qua em dọn bếp xong là lên phòng ngủ luôn, không có lên tầng năm đâu.''

Từ dưới cầu thang, tiếng Ryuuhou vọng lên, rõ ràng là đang nghe từ tầng năm: ``Em cũng không lên đâu nhé! Đêm qua em còn xử lý báo cáo đến khuya, không có bước chân lên tầng năm đâu!''

Từ phòng khách, Kson cũng đáp lại, em ấy vừa về từ kho tối qua nên vẫn đang sắp xếp đồ đạc: ``Em còn ở kho đến tận tối mịt mới về, không thể vào phòng anh được đâu ạ.''

Delutaya, người vừa ra khỏi phòng thu của mình, cũng xác nhận: ``Tôi đã khóa phòng thu từ trước khi Kuon-kun ngủ say, cả đêm không ra khỏi phòng nên không phải tôi đâu.''

Còn hai đứa nhỏ Naomi và Koneko thì vẫn đang ngủ say trên tầng tám, chẳng thể nào xuống tầng năm giữa đêm được. Tất cả mọi người trong nhà đều phủ nhận, không ai nhận là người đã mang nước và thuốc vào phòng cho tôi.

Tôi vô thức ngẩng đầu nhìn lên cầu thang đá dẫn đến tầng tám, nơi Mikeneko đang ở. Game trên đồng hồ im lặng, không nói thêm gì. Tôi cũng không hỏi thêm ai nữa, có một cảm giác nhẹ nhàng lan trong lòng, tôi dường như đã đoán ra ai là người đã làm việc đó.

\ct

Buổi sáng trôi qua một cách chậm rãi, bình yên như những vệt nắng di chuyển chậm trên sàn nhà. Tôi vẫn tiếp tục nằm nghỉ ngơi theo đúng lời dặn của bác sĩ ở phòng khám, không dám vận động nhiều. Nhiệt độ cơ thể của tôi giảm dần qua mỗi lần đo, từ hơn 39 độ xuống còn 37.8 độ, từng chút một trở lại mức bình thường.

Suzuno mang lên một tô cháo nóng hổi, thơm phức mùi hành tây và gừng, đặt cạnh bàn. Naomi sau đó cũng chạy lên, đặt một bức tranh vẽ màu rực rỡ cạnh chiếc cốc nước - đó là hình ảnh một con gấu trúc nằm nghỉ ngơi dưới cây cổ thụ, em bé viết dòng chữ nhỏ ``Onii-chan mau khỏe'' ở góc tranh.

Kaigou cũng không còn kiểm tra trán tôi mỗi mười lăm phút như hôm qua, ít nhất là là em ấy cố gắng che giấu việc đó, không để tôi thấy vẫn luôn lo lắng.

Ryuuhou gửi cho tôi một danh sách chi tiết những công việc đã được chuyển giao, tất cả các cuộc họp và đơn hàng đều được xử lý ổn thỏa, không có gì bị trì hoãn.

Kson nhắn tin từ kho, báo cáo ca làm việc hôm nay cũng kết thúc tốt đẹp, mọi thứ ở kho đều hoạt động trơn tru. Delutaya còn để một bình trà hoa cúc nhẹ bên ngoài cửa phòng, hương thơm dịu nhẹ bay vào mỗi khi có gió thổi qua.

Cả nhà đều bận rộn với phần việc của mình, và không ai biết ai đã là người mang cốc nước đêm qua. Câu hỏi đó lặng lẽ trôi qua, không ai nhắc lại nữa. Và tôi cũng không cần phải biết chắc chắn, sự ấm áp từ việc có người quan tâm thầm lặng đã đủ làm lòng tôi nhẹ nhõm.

Đến chiều, cơn sốt gần như đã lui hẳn. Tôi vẫn còn cảm thấy mệt mỏi, nhưng đã có thể ngồi dậy trên giường mà không thấy căn phòng quay vòng, không còn cảm giác choáng váng mỗi khi cử động. Kaigou dựa vào khung cửa, khoanh tay trước ngực với vẻ mặt vẫn cố gắng nghiêm nghị: ``Anh thấy đỡ hơn chưa? Đừng có tưởng đỡ là được xuống làm việc ngay đâu nhé.''

Tôi cười nhẹ, gật đầu: ``Nhiều rồi. Anh chỉ định ngồi một lúc thôi, không có định đi đâu.''

``Ngồi trên giường thôi, đừng có đứng dậy lung tung,'' em nhắc lại, giọng vẫn không mềm đi.

Tôi bật cười khẽ, nhưng cơn cười kéo theo một tiếng ho nhỏ bật ra khỏi cổ họng. Ngay lập tức, Kaigou bước nhanh đến cạnh giường, tay đặt lên trán tôi kiểm tra nhiệt độ, giọng có chút trách móc: ``Thấy chưa? Vẫn còn ho mà. Nằm xuống ngay đi, đừng cố gắng quá.''

Tôi làm theo lời em ấy, nằm lại trên giường, kéo tấm chăn ấm lên đến ngực. Mặc dù luôn nói những lời mạnh mẽ, nhưng sự lo lắng trong mắt em ấy luôn rõ ràng, và tôi cảm thấy ấm áp vì có cả nhà ở bên cạnh chăm sóc.

\ct

Khi màn đêm buông xuống, bóng tối dần bao phủ từng góc căn nhà, tiếng cười nói thoảng vọng lên từ phòng khách nơi mọi người đã tập trung dùng bữa tối. Âm thanh của thìa đũa va chạm nhẹ và những câu chuyện nhỏ len lỏi qua khe cửa, nhưng tôi vẫn chọn ở lại trong phòng, vừa ăn xong bát cháo nóng hổi Suzuno mang lên và uống đủ cốc nước ấm để cổ họng bớt đi cơn khô rát.

Chẳng mấy chốc, tiếng bước chân nhẹ đến trước cửa, Kaigou ghé vào đây là lượt thăm cuối cùng trước khi em ấy lên phòng ngủ. Giọng em vẫn giữ vẻ nghiêm nghị như mọi khi, nhưng ẩn dưới đó là nỗi lo chưa hề tan: ``Mai anh vẫn phải nghỉ nguyên buổi sáng. Đừng cãi lời em.''

Tôi thở dài một tiếng, gật đầu thừa nhận: ``Anh biết rồi, không cãi đâu.''

``Game sẽ khóa toàn bộ lịch trình của anh, không để anh động vào công việc đâu,'' em ấy nhấn mạnh thêm.

Ngay lập tức, giọng điện tử đều đặn của Game vang lên từ chiếc đồng hồ thông minh trên bàn, xác nhận lời của em ấy: ``Tôi đã khóa toàn bộ lịch thành công.''

Sau khi nghe vậy, Kaigou bước lại gần giường, kéo nhẹ tấm chăn mỏng lên đến ngang vai tôi, che chắn cho tôi khỏi hơi lạnh của đêm sớm. ``Ngủ sớm đi nhé.''

``Em cũng nên lên nghỉ ngơi sớm đi, ngày hôm nay cũng vất vả lắm rồi,'' tôi đáp lại, nhìn bóng dáng em ấy vẫn đứng yên ở mép giường.

``Em biết rồi,'' Kaigou nói, nhưng đôi chân vẫn chưa bước đi, em ấy đứng lặng đó thêm vài giây, như thể còn muốn kiểm tra lại một lần nữa tình trạng của tôi.

Tôi nhắm chặt mắt, giả vờ đã chìm vào giấc ngủ, không muốn em ấy phải bận tâm thêm lúc nào nữa. Cuối cùng, Kaigou mới rón rén bước ra khỏi phòng, tiếng cửa khép lại khe khẽ như muốn không làm phiền tôi.

Cả căn phòng lập tức chìm vào im lặng, chỉ còn tiếng tích tắc nhỏ của đồng hồ treo tường. Tôi mở mắt trong bóng tối, đầu óc không thể ngừng nghĩ về chiếc cốc nước ấm và viên thuốc hạ sốt mà người lạ đặt cạnh giường tối hôm trước. Trong trí nhớ, tôi phác họa lại hình ảnh cầu thang đá dẫn lên các tầng trên, suy nghĩ mông lung về người đã lặng lẽ làm việc đó mà không để lại dấu vết. Cả nhà đều phủ nhận, chẳng ai nhận là mình đã vào phòng. Một nụ cười nở trên môi tôi, và chẳng mấy chốc, cơn mệt mỏi kéo tôi chìm vào giấc ngủ say.

\ct

Sáng hôm sau, khi ánh nắng đầu tiên len lỏi qua khe cửa, nhiệt độ cơ thể của tôi đã gần như quay về mức bình thường. Cơn sốt nóng bức đã hoàn toàn lui đi, chỉ còn lại cảm giác đầu hơi nặng nề, mệt mỏi như vừa trải qua một cuộc chạy đua dài.

Game đọc lại kết quả đo nhiệt độ từ chiếc nhiệt kế điện tử, giọng vẫn đều đặn nhưng có chút quan tâm: ``Nhiệt độ cơ thể đã ổn định. Nếu vẫn còn cảm thấy mệt, thì cứ đừng ngần ngại mà nghỉ ngơi.''

Tại cửa phòng, Kaigou đứng đó, vẻ mặt nhẹ nhõm rõ rệt nhưng em ấy vẫn cố gắng giữ cho giọng mình nghiêm nghị, như sợ tôi lại đòi làm việc: ``Dù anh đã đỡ hơn nhiều, cũng không được ngồi trước máy tính cả ngày đâu.''

Tôi cười khẽ, cố gắng thương lượng: ``Anh chỉ định mở máy xem một vài email quan trọng thôi, không làm gì nhiều.''

``Không được,'' Kaigou lập tức phản đối không chút do dự.

``Chỉ một email thôi mà, em quá đáng quá,'' tôi vẫn cố gắng thuyết phục.

``Cũng không được,'' em ấy lắc đầu kiên quyết.

Ngay lúc đó, tiếng Ryuuhou vọng lên từ cầu thang, như thể đã chờ đợi cuộc đối thoại này từ lâu: ``Em đã trả lời hết tất cả những email gấp rồi anh ạ. Anh không cần phải mở máy đâu, mọi việc đã ổn cả rồi.''

Tôi quay sang nhìn Kaigou, em ấy nhún vai một cách vô tội, như thể chuyện đó chẳng liên quan gì đến mình, khiến tôi bật cười. Chẳng mấy chốc, Suzuno mang bữa sáng nóng hổi lên phòng, còn Naomi chạy theo sau, đặt một bức tranh vẽ mới màu rực rỡ cạnh chiếc cốc sứ cũ.

Tôi nhìn quanh căn phòng, từng người một đều đã lo liệu hết phần việc của mình, chẳng để lại gì cho tôi phải bận tâm. Tôi không cần phải cố gắng đứng dậy gánh vác mọi thứ ngay lập tức, mọi người đã chở che cho tôi quá tốt. Tôi nằm lại trên giường, lần này không còn ý định thương lượng gì nữa, chỉ yên tâm tận hưởng khoảng thời gian nghỉ ngơi này.

\ct

Đến chiều tà, khi ánh nắng đã chuyển sang màu cam ấm áp, tôi đã có thể tự mình bước xuống phòng khách ngồi nghỉ một lúc. Cả nhà dường như đều thở phào nhẹ nhõm khi thấy tôi đã khỏe lại nhiều.

Kaigou ngồi cạnh tôi, kể lại rằng em ấy đã thay tôi xử lý vài cuộc gọi khách hàng khó tính, mọi thứ đều diễn ra thuận lợi. Suzuno từ bếp gọi vọng ra, nói rằng trong nồi vẫn còn cháo nóng, nếu tôi thấy đói cứ nói để em ấy mang ra. Ryuuhou cũng cập nhật rằng tất cả các lịch họp quan trọng đã được chuyển giao cho người khác, không có cuộc nào bị trì hoãn.

Kson nhắn tin từ kho, báo cáo ca làm việc hôm nay cũng kết thúc tốt đẹp, mọi đơn hàng đều được giao đúng giờ. Delutaya còn mang ra một bình trà hoa cúc mới ủ, hỏi tôi có muốn rót thêm một cốc không để cơ thể thêm ấm. Cuối cùng, Naomi chạy đến đặt vào tay tôi một bức tranh mới, vẽ hình hai người đang cười nói dưới bóng cây, nét vẽ nguệch ngoạc nhưng tràn đầy tình cảm.

Tôi lần lượt cảm ơn từng người, lòng ấm áp vô cùng. Không ai trong nhà nhắc lại câu hỏi về người đã mang cốc nước và thuốc lên phòng đêm đó. Tôi cũng không bao giờ hỏi lại nữa. Chỉ có chiếc cốc sứ cũ được rửa sạch, nằm gọn trên khay đựng chén đĩa cạnh cửa.

Trong đầu tôi lại hiện lên hình ảnh một người lặng lẽ bước qua hành lang tối mịn, đặt cốc nước cạnh giường rồi lặng lẽ rời đi, không muốn ai biết đến việc mình đã làm. Ánh mắt tôi vô thức hướng lên cầu thang dẫn đến tầng tám, nơi Mikeneko đang ở. Tôi không cần phải hỏi để biết chắc chắn sự thật. Trong lòng, tôi đã hiểu rõ tất cả.

%==========================================TRUYỆN PHỤ=========================================
\omk

Đêm buông xuống hoàn toàn, cả nhà Hikawa quây quần trong phòng khách ấm áp sau bữa tối thịnh soạn. Mùi hương của trà hoa cúc vẫn bay nhẹ trong không khí, ánh đèn vàng dịu dàng làm tan đi cái lạnh của đêm sớm. Game đợi đến khi tất cả mọi người đều ngồi yên trên ghế dài, chưa kịp lan ra về phòng của mình, thì mới bật chiếc máy chiếu nhỏ đã được chuẩn bị từ trước.

Ánh sáng trắng nhè nhẹ hiện ra trên bức tường trắng, chiếu rõ đoạn ghi hình của camera hành lang tầng năm, quay lại từ những ngày tôi bị ốm. Mọi người lập tức dồn mắt về phía màn hình, tò mò xem có gì mà Game lại đột nhiên chiếu lên.

Trên màn hình, Kaigou hiện ra trước cửa phòng, dáng đi lo lắng như đang đi qua đi lại không ngừng. Cứ đúng mười lăm phút trôi qua, em ấy lại nhẹ nhàng đẩy cửa phòng bước vào, chỉ ở lại vài giây để đặt lòng bàn tay lên trán tôi kiểm tra cơn sốt, rồi lại lặng lẽ bước ra, đóng cửa thật chậm để không làm phiền. Đến lượt thứ tư em ấh vào phòng, tay còn ôm theo một ấm nước gừng lớn còn bốc hơi nóng, miệng lẩm bẩm gì đó với chính mình trước khi đặt ấm nước lên bàn cạnh giường.

``\textbf{Game! Tắt ngay!}'' Kaigou bật dậy khỏi ghế bành, mặt bừng đỏ lên khi nhận ra mình chính là nhân vật chính trong đoạn ghi hình. Em vừa hét vừa lao về phía chiếc đồng hồ thông minh của Game đang đặt trên bàn trà, muốn ngăn chặn cái màn chiếu xấu hổ này càng sớm càng tốt.

Nhưng không những không tắt, ông ta còn tua lại đúng đoạn Kaigou đứng cạnh giường tôi, tay vừa rút ra sau khi chạm trán, tự nhủ nhỏ với mình rằng phải bình tĩnh, không được hoảng loạn chỉ vì cơn sốt thông thường. Giọng điện tử đều đặn của vang lên trong phòng khách, tràn đầy vẻ hài lòng: ``Tôi đã lưu lại hồ sơ chăm sóc người nhà cực kỳ cẩn thận. Đây là dữ liệu quý giá.''

``\textbf{Ông xóa nó đi ngay lập tức!}'' Kaigou hét lên, tay đã chạm vào mép chiếc đồng hồ, nhưng tôi ngồi cạnh đã nhanh tay nhấc nó lên, giơ ra xa khỏi tầm với của em.

``Tôi có thể tạo thêm một bản sao dự phòng để lưu vào đám mây, đề phòng mất dữ liệu,'' Game tiếp tục giỡn, không chịu dừng đoạn ghi hình.

``Game!'' Kaigou càng tức giận, nhảy lên cố gắng với lấy chiếc đồng hồ trên tay tôi. ``Aniki, đưa cái đồng hồ đó cho em đi! Ông ấy đang phá bĩnh quá!''

Tôi cười lớn, giơ chiếc đồng hồ lên cao hơn nữa, để Kaigou không thể với tới. ``Bình tĩnh nào Kaigou, anh nghĩ đoạn này khá đáng yêu mà. Ai lại không có người em gái lo lắng chăm sóc mình thế này đâu.''

``Em còn không thèm lo! Em chỉ là làm bổn phận thôi!'' em ấy vùng vằng, mặt càng đỏ hơn khi nghe mọi người trong phòng cùng cười lên. Ryuuhou cười đến mức phải ôm bụng, lăn ra trên ghế dài, không kiềm được tiếng cười to. Suzuno cũng che miệng cười nhẹ, nhìn Kaigou với ánh mắt trìu mến. Ngay cả hai đứa nhỏ Naomi và Koneko cũng vỗ tay cười rộ, không hiểu rõ chuyện gì nhưng cũng thấy vui theo không khí.

Mikeneko ngồi ở góc phòng, một mình bên cửa sổ, cũng khóe miệng nở một nụ cười nhỏ. Ánh mắt cô lướt qua tôi đang cười lớn, rồi lại quay về Kaigou đang tức giận đòi lấy lại chiếc đồng hồ. 

Trong căn phòng ấm áp đầy tiếng cười, mọi góc tối của căn nhà Hikawa đều được sưởi ấm bởi tình cảm của mọi người dành cho nhau. Câu chuyện nhỏ về cốc nước ấm đêm đó vẫn nằm gọn trong lòng hai người, không cần phải nói ra, vì tình cảm thầm lặng ấy đã hòa vào không khí ấm áp của cả gia đình, trở thành một phần nhỏ bé nhưng đáng quý trong những ngày bình yên này.

Cuối cùng tôi cũng nhượng bộ, đặt chiếc đồng hồ xuống bàn, nhưng không quên nháy mắt với Game trước khi để Kaigou giật lấy. Ông ta lập tức tắt máy chiếu, nhưng tiếng cười vẫn còn vang vọng trong phòng khách rất lâu sau đó, làm cho đêm đông lạnh giá trở nên ấm áp hơn bao giờ hết.

\end{document}